% concatenated from main.tex
\documentclass[12pt]{amsart}

\usepackage[T1]{fontenc}
\usepackage{amsmath,amssymb,amsthm}
\usepackage{aliascnt,enumerate,shuffle}
\usepackage[margin=24mm]{geometry}
\usepackage{xcolor}
\definecolor{mylinkcolor}{RGB}{16, 156, 81}
\definecolor{mycitecolor}{RGB}{20, 80, 140}
\usepackage{hyperref}
\usepackage[nameinlink]{cleveref}
\hypersetup{
    setpagesize=false,
    bookmarksnumbered=true,
    bookmarksopen=true,
    colorlinks=true,
    linkcolor=mylinkcolor,
    citecolor=mycitecolor,
}

\numberwithin{equation}{section}
\allowdisplaybreaks[2]

\theoremstyle{plain}
\newtheorem{thm}{Theorem}[section]
\crefname{thm}{Theorem}{Theorems}

\newaliascnt{lem}{thm}
\newtheorem{lem}[lem]{Lemma}
\aliascntresetthe{lem}
\crefname{lem}{Lemma}{Lemmas}

\newaliascnt{prop}{thm}
\newtheorem{prop}[prop]{Proposition}
\aliascntresetthe{prop}
\crefname{prop}{Proposition}{Propositions}

\newaliascnt{cor}{thm}
\newtheorem{cor}[cor]{Corollary}
\aliascntresetthe{cor}
\crefname{cor}{Corollary}{Corollaries}

\newaliascnt{conj}{thm}
\newtheorem{conj}[conj]{Conjecture}
\aliascntresetthe{conj}
\crefname{conj}{Conjecture}{Conjectures}

\newaliascnt{dfn}{thm}
\newtheorem{dfn}[dfn]{Definition}
\aliascntresetthe{dfn}
\crefname{dfn}{Definition}{Definitions}

\newaliascnt{rem}{thm}
\newtheorem{rem}[rem]{Remark}
\aliascntresetthe{rem}
\crefname{rem}{Remark}{Remarks}


\newaliascnt{ex}{thm}
\newtheorem{ex}[ex]{Example}
\aliascntresetthe{ex}
\crefname{ex}{Example}{Examples}

\newtheorem{mainthm}{Main Theorem}
\renewcommand{\themainthm}{\Alph{mainthm}}
\crefname{mainthm}{Main Theorem}{Main Theorems}
\newaliascnt{mainconj}{mainthm}
\newtheorem{mainconj}[mainconj]{Main Conjecture}
\aliascntresetthe{mainconj}
\crefname{mainconj}{Main Conjecture}{Main Conjectures}

\crefname{section}{Section}{Sections}
\crefname{subsection}{Section}{Sections}

\newcommand{\odd}{\mathrm{odd}}
\newcommand{\even}{\mathrm{even}}
\newcommand{\QQ}{\mathbb{Q}}
\newcommand{\ZZ}{\mathbb{Z}}
\DeclareMathOperator{\Span}{span}
\DeclareMathOperator{\SL}{SL}
\DeclareMathOperator{\Fil}{Fil}
\DeclareMathOperator{\wt}{wt}
\DeclareMathOperator{\dep}{dep}

\newcommand{\h}{\mathfrak H}
\newcommand{\A}{\mathcal A}
\newcommand{\ZA}{{\mathcal Z_{\mathcal A}}}
\newcommand{\Zf}{{\mathcal Z^{\mathrm f}}}
\newcommand{\ZAf}{{\mathcal Z_{\mathcal A}^{\mathrm f}}}
\newcommand{\ZAw}[1]{{\mathcal Z_{\mathcal A,#1}}}
\newcommand{\ZAfw}[1]{{\mathcal Z_{\mathcal A,#1}^{\mathrm f}}}
\newcommand{\Zbarf}{{\overline{\mathcal Z}^{\mathrm f}}}
\newcommand{\za}[1]{\zeta_{\mathcal A}(#1)}
\newcommand{\zs}[1]{\zeta_{\mathcal S}(#1)}
\newcommand{\zf}[1]{\zeta^{\mathrm f}(#1)}
\newcommand{\zaf}[1]{\zeta_{\mathcal A}^{\mathrm f}(#1)}
\newcommand{\zsf}[1]{\zeta_{\mathcal S}^{\mathrm f}(#1)}
\newcommand{\zzf}[1]{\mathfrak z^{\mathrm f}(#1)}
\newcommand{\Dk}{\mathcal D_k}
\newcommand{\Dtk}{\widetilde{\mathcal D}_k}
\newcommand{\Dbk}{\overline{\mathcal D}_k}
\newcommand{\emptyword}{\mathbf 1}

\title{Formal finite multiple zeta values}

\author{Henrik Bachmann}
\address{Graduate School of Mathematics, Nagoya University, Nagoya, Japan.}
\email{henrik.bachmann@math.nagoya-u.ac.jp}

\author{Risan}
\address{Nagoya, Japan.}
\email{ptrrsn.1@gmail.com}

\date{}

\subjclass[2020]{Primary 11M32; Secondary 11F11}
\keywords{finite multiple zeta values, symmetric multiple zeta values, formal multiple zeta values, Kaneko-Zagier conjecture, period polynomials}

%================================================
%                 Document Body
%================================================
\begin{document}

\begin{abstract}
We study the algebra of formal finite multiple zeta values by imposing the
stuffle and linear shuffle relations of Kaneko and Zagier. Our first results
are a surjective homomorphism to the algebra of formal symmetric multiple
zeta values and a parity reduction in depth at most four. In even weight we
construct a quotient of the formal double zeta space which surjects onto the
space of finite multiple zeta values of depth at most four, and we attach to
every even period polynomial an explicit relation among the values
$\zeta_{\mathcal A}(2a,1,2b,1)$.
\end{abstract}

\maketitle

\section{Introduction}

For an index $(k_1,\ldots,k_r) \in \ZZ_{\geq 1}^r$ with $k_1\geq2$, the \emph{multiple zeta
value} is defined by
\begin{align}\label{eq:intro-mzv}
 \zeta(k_1,\ldots,k_r)
 =\sum_{m_1>\cdots>m_r>0}
 \frac{1}{m_1^{k_1}\cdots m_r^{k_r}}.
\end{align}
We write $\mathcal Z$ for the $\QQ$-algebra generated by these values.
Multiple zeta values satisfy the stuffle and shuffle product relations. Their
regularized combination gives the extended double shuffle relations of Ihara,
Kaneko and Zagier~\cite{IKZ}, which are conjectured to generate all algebraic
relations among multiple zeta values.

Kaneko and Zagier introduced two finite counterparts of
\eqref{eq:intro-mzv}~\cite{KZ}. The first are the finite multiple zeta values
\begin{align}\label{eq:intro-fmzv}
 \za{k_1,\ldots,k_r}
 =\left(\sum_{p>m_1>\cdots>m_r>0}
 \frac{1}{m_1^{k_1}\cdots m_r^{k_r}}\bmod p\right)_p
 \in\A,
\end{align}
where
\[
 \A=\prod_{p\text{ prime}}\ZZ/p\ZZ\bigg/\bigoplus_{p\text{ prime}}\ZZ/p\ZZ.
\]
The second are the symmetric multiple zeta values
$\zs{k_1,\ldots,k_r}\in\mathcal Z/\zeta(2)\mathcal Z$, defined as alternating
sums of products of regularized multiple zeta values. Although the two
definitions are very different, both families satisfy the stuffle product and
the same linear shuffle relations. Kaneko and Zagier conjecture that these
relations give all relations among finite multiple zeta values and that the
assignment
\begin{align}\label{eq:intro-kz}
 \za{k_1,\ldots,k_r}\longmapsto\zs{k_1,\ldots,k_r}
\end{align}
induces an isomorphism
$\ZA\simeq\mathcal Z/\zeta(2)\mathcal Z$. It is not known that
\eqref{eq:intro-kz} is well defined.

Motivated by these conjectures, we define the algebra
$\ZAf$ of \emph{formal finite multiple zeta values} by imposing precisely the
stuffle and linear shuffle relations, and we denote the class of an index by
$\zaf{k_1,\ldots,k_r}$. Starting from the algebra $\Zf$ of formal multiple
zeta values, we also define formal symmetric multiple zeta values
\[
 \zsf{k_1,\ldots,k_r}\in
 \Zbarf:=\Zf/\zf{2}\Zf.
\]
Our first result is a formal version of the Kaneko--Zagier correspondence.

\begin{mainthm}\label{main:formal-kz}
There is a surjective homomorphism of graded $\QQ$-algebras
\[
 \begin{aligned}
 \varphi:\ZAf&\longrightarrow\Zbarf,\\
 \zaf{k_1,\ldots,k_r}&\longmapsto\zsf{k_1,\ldots,k_r}.
 \end{aligned}
\]
The classical finite and symmetric multiple zeta values are quotient
realizations of $\ZAf$.
\end{mainthm}

The surjectivity is the formal version of Yasuda's theorem. Yasuda notes that
his proof uses only regularized double shuffle relations and therefore applies
to any realization satisfying them, see the introduction and the proof of the
main theorem of~\cite{Yasuda}. Applying
that argument to the universal formal algebra $\Zf$, we obtain the result above.
Injectivity remains open and is the formal analogue of the Kaneko--Zagier
conjecture.

Our second result concerns the depth filtration. The Kaneko--Zagier
correspondence is expected to decrease depth by one. Together with the parity
theorem for multiple zeta values~\cite{Tsumura,Panzer}, this suggests that a
formal finite value of weight $k$ and depth $d$ should reduce to lower depth
when $k\equiv d\pmod2$. The corresponding statement is known for symmetric
multiple zeta values in every depth~\cite[Corollary to Theorem~4]{KZ}, but it
does not imply the result in the universal formal algebra or for finite
multiple zeta values. Denote by $\Fil_d\ZAfw{k}$ the span of formal finite multiple zeta values of
weight $k$ and depth $\leq d$, and by $\Fil_d\ZAw{k}$ the analogous span of
finite multiple zeta values.

\begin{mainthm}\label{main:parity}
Let $k,d\geq1$ with $k\equiv d\pmod2$. If $d\leq4$, then
\[
 \Fil_d\ZAfw{k}\subseteq\Fil_{d-1}\ZAfw{k}.
\]
\end{mainthm}

The proof first reduces depth four to products of two odd-weight formal single
values and then eliminates these products by one generating-series form of the
depth $1+3$ linear shuffle relation. It works entirely in $\ZAf$ and hence
specializes simultaneously to finite and symmetric multiple zeta values.

The remaining part of the paper concerns period polynomials. For even
$k\geq4$, let $W_k^{\even}$ be the space of homogeneous polynomials of degree
$k-2$ satisfying
\begin{align}\label{eq:even-period-polynomials}
 p(X,Y)=p(X+Y,Y)-p(X+Y,X).
\end{align}
Every element is antisymmetric and contains only even powers of $X$ and $Y$,
and
\[
 \dim_{\QQ}W_k^{\even}=\dim_{\QQ}\mathcal S_k+1,
\]
where $\mathcal S_k$ is the space of cusp forms of weight $k$ for
$\SL_2(\ZZ)$. The line $\QQ(X^{k-2}-Y^{k-2})$ is the noncuspidal part.
Gangl, Kaneko and Zagier showed that even period polynomials give relations
among double zeta values~\cite{GKZ}, see also~\cite{BS,Tasaka} for further
period-polynomial relations. Kaneko and Zagier conjecture that the
finite multiple zeta values
\[
 \za{2a,1,2b,1},\qquad a,b\geq1,\quad 2a+2b+2=k,
\]
span the depth-four part in weight $k$ and satisfy exactly
$\dim_{\QQ}\mathcal S_k$ independent relations~\cite{KZ}.
On the symmetric side, Hara, Sakugawa and Tasaka recently expressed
every cusp form with rational Fourier coefficients as a rational
linear combination of symmetric triple Eisenstein series, whose
constant terms give modular relations among symmetric triple zeta
values~\cite{HST}.

To explain the relations among the values $\za{2a,1,2b,1}$, let $\Dk$
be the formal double zeta space of
Gangl, Kaneko and Zagier and define
\[
 \Dtk=\Dk\big/\left(
 \QQ Z_k+\QQ Z_{1,k-1}
 +\Span_{\QQ}\{P_{r,s}:r,s\text{ even},\ r+s=k\}\right).
\]
Its dimension is
\[
 \dim_{\QQ}\Dtk=\frac{k}{2}-2-\dim_{\QQ}\mathcal S_k.
\]
For odd $r,s\geq3$ with $r+s=k$, put
\begin{align}\label{eq:H-source-element}
 H_{r,s}=\frac12\Bigg(&
 (rs-r+1)Z_{r,s}+(rs+s-1)Z_{s,r}
 +(s+1)Z_{s+1,r-1}-(r+1)Z_{r+1,s-1}\notag\\
 &+\sum_{\substack{3\leq\ell\leq k-3\\\ell\ \mathrm{odd}}}
 \left(\binom{k-\ell}{s}-\binom{k-\ell}{r}\right)
 P_{\ell,k-\ell}\Bigg)\in\Dtk.
\end{align}
These elements can be seen as the counterparts in $\Dtk$ of the values
$\zeta_{\mathcal A}(r-1,1,s-1,1)$.

\begin{mainthm}\label{main:double-zeta-surjection}
Let $k\geq4$ be even.
\begin{enumerate}[(i)]
\item The elements $H_{r,s}$, where $r,s\geq3$ are odd and $r+s=k$, span
$\Dtk$.
\item Assume that \cref{conj:beta-identities} holds in weight $k$. Then 
\[
 H_{r,s}\longmapsto\zaf{r-1,1,s-1,1}
\]
defines a surjective map
$\beta_k^{\mathrm f}:\Dtk\to\Fil_4\ZAfw{k}$.
\end{enumerate}
\end{mainthm}

For classical finite multiple zeta values, the depth-three identity needed in
part~(ii) follows from Hoffman's duality together with formulas of
Kina~\cite{HoffmanFinite,Kina}. This gives an unconditional map in every even
weight. The coordinates realizing this map agree, after reversal, with
depth-three elements introduced by Kina, who proved their shuffle relation
and that they span the depth-three values of even weight~\cite{Kina}. Our new
input is the complementary stuffle relation for finite multiple zeta values,
which turns these coordinates into a realization of the reduced formal double
zeta space, see \cref{rem:kina-comparison}.

\begin{mainthm}\label{main:classical-beta}
For every even $k\geq4$, 
\[
 H_{r,s}\longmapsto\za{r-1,1,s-1,1}
\]
defines a surjective map
$\beta_k:\Dtk\to\Fil_4\ZAw{k}$. Consequently,
\[
 \Fil_4\ZAw{k}
 =\Span_{\QQ}\{\za{r-1,1,s-1,1}:r,s\geq3\text{ odd},\ r+s=k\},
\]
and these values satisfy at least $\dim_{\QQ}\mathcal S_k$ independent
linear relations.
\end{mainthm}

We expect the two maps to identify the reduced formal double zeta space with
the corresponding depth part.

\begin{mainconj}\label{main:double-zeta}
For every even $k\geq4$, the rule defining $\beta_k^{\mathrm f}$ is well
defined, and the maps $\beta_k^{\mathrm f}$ and $\beta_k$ are isomorphisms.
\end{mainconj}

The period-polynomial theorem of~\cite{GKZ} gives all relations among the
classes $Z_{r,k-r}$ with $r,k-r\geq3$ odd in $\Dtk$. We rewrite the $H_{r,s}$ in
these classes by an explicit matrix with binomial and Bernoulli number entries.
Cramer's rule then assigns to every even period polynomial an explicit relation
among the $H_{r,s}$, and applying $\beta_k$ gives an unconditional relation
among finite multiple zeta values. Whenever this matrix is invertible, which
we verified by exact computation for every even $k\leq100$
(\cref{conj:detA}), these relations
are nontrivial for all cuspidal period polynomials and exhaust the relations
among the $H_{r,s}$. For example, for the period polynomial
$X^2Y^2(X^2-Y^2)^3$ of weight $12$, our construction recovers the relation
\begin{align}\label{eq:intro-weight-twelve-relation}
 16\za{2,1,8,1}+9\za{4,1,6,1}
 +18\za{6,1,4,1}-2\za{8,1,2,1}=0,
\end{align}
which was observed by Kaneko and Zagier~\cite{KZ}.
If $\beta_k$ is bijective, these are all relations among the values
$\za{2a,1,2b,1}$, and the conjecture of Kaneko and Zagier follows.

The organization of this paper is as follows. In \cref{sec:prelim} we recall finite,
symmetric and formal multiple zeta values. In \cref{sec:formal-finite} we
define $\ZAf$ and prove \cref{main:formal-kz}. The parity theorem is proved in
\cref{sec:parity}. In \cref{sec:double-zeta} we study $\Dtk$, construct the
maps $\beta_k^{\mathrm f}$ and $\beta_k$, and derive the period-polynomial
relations.

\Cref{main:formal-kz,main:parity} are contained in the master's thesis
\cite{Risan} of the second author, written under the supervision of the first
author. 

\subsection*{Acknowledgements}
We thank Masanobu Kaneko for fruitful discussions on this topic. This
project was partially supported by JSPS KAKENHI Grants 23K03030 and
26K22254.

\section{Finite, symmetric and formal multiple zeta values}\label{sec:prelim}

Throughout the paper an \emph{index} is a tuple $\boldsymbol{k}=(k_1,\ldots,k_r)$ of positive integers. Its weight and depth are
\[
 \wt(\boldsymbol{k})=k_1+\cdots+k_r,
 \qquad \dep(\boldsymbol{k})=r.
\]
We use decreasing summation indices as in \eqref{eq:intro-fmzv}. Note that Kaneko and Zagier use increasing summation indices in~\cite{KZ}, and the two conventions are related by reversal.

\subsection{The two products}

Let $\h=\QQ\langle x,y\rangle$ and
\[
 \h^0=\QQ+x\h y,
 \qquad
 \h^1=\QQ+\h y.
\]
We put $z_k=x^{k-1}y$ for $k\geq1$, i.e. the words in $\h^1$ are $\emptyword$ and $z_{k_1}\cdots z_{k_r}$. The \emph{stuffle product} is the bilinear product on $\h^1$ determined by
\begin{align}\label{eq:stuffle}
 \emptyword*w=w*\emptyword=w,
 \qquad
 z_aw*z_bv=z_a(w*z_bv)+z_b(z_aw*v)+z_{a+b}(w*v).
\end{align}
The \emph{shuffle product} on $\h$ is determined by
\begin{align}\label{eq:shuffle}
 \emptyword\shuffle w=w\shuffle\emptyword=w,
 \qquad
 au\shuffle bv=a(u\shuffle bv)+b(au\shuffle v)
\end{align}
for $a,b\in\{x,y\}$. We denote the corresponding commutative algebras by $\h^1_*$ and $\h^1_\shuffle$.

Let
\[
 R(z_{k_1}\cdots z_{k_r})=z_{k_r}\cdots z_{k_1}
\]
be reversal. The deconcatenation coproduct
\[
 \Delta(w)=\sum_{uv=w}u\otimes v
\]
makes $\h^1_*$ a Hopf algebra~\cite{HoffmanIhara}. If $\Sigma$ replaces any collection of separators in a reversed word by contractions $z_a\diamond z_b=z_{a+b}$ and $\epsilon_{\dep}(w)=(-1)^{\dep(w)}w$, its antipode is
\begin{align}\label{eq:antipode}
 S_*=\Sigma\circ\epsilon_{\dep}\circ R.
\end{align}
We will use the antipode identity
\begin{align}\label{eq:antipode-identity}
 \sum_{uv=w}S_*(u)*v=0
\end{align}
for every nonempty word $w$.

\subsection{Finite and symmetric values}

The map $Z_{\mathcal A}:\h^1\to\A$ defined by \eqref{eq:intro-fmzv} is an algebra homomorphism for the stuffle product. It also satisfies the linear shuffle relation of Kaneko and Zagier: for all words $w,v\in\h^1$,
\begin{align}\label{eq:finite-linear-shuffle}
 Z_{\mathcal A}(w\shuffle v)=(-1)^{\wt(w)}Z_{\mathcal A}(R(w)v).
\end{align}
In particular,
\begin{align}\label{eq:finite-reversal}
 \za{k_r,\ldots,k_1}=(-1)^{k_1+\cdots+k_r}\za{k_1,\ldots,k_r}.
\end{align}
We write $\ZA\subset\A$ for the $\QQ$-algebra spanned by all finite multiple zeta values.

To recall the second finite version, let $\zeta^{*,T}$ and $\zeta^{\shuffle,T}$ denote the stuffle and shuffle regularizations of multiple zeta values. For $\bullet\in\{*,\shuffle\}$ set
\begin{align}\label{eq:symmetric-classical}
 \zeta_{\mathcal S}^{\bullet}(k_1,\ldots,k_r)
 =\sum_{j=0}^r(-1)^{k_1+\cdots+k_j}
 \zeta^{\bullet,T}(k_j,\ldots,k_1)
 \zeta^{\bullet,T}(k_{j+1},\ldots,k_r),
\end{align}
where the value attached to the empty index is~$1$. This expression is independent of $T$, and the two choices of regularization are congruent modulo $\zeta(2)\mathcal Z$. Their common class is $\zs{k_1,\ldots,k_r}$. Symmetric multiple zeta values satisfy the stuffle product and \eqref{eq:finite-linear-shuffle}, see~\cite{KZ,Kaneko}.

\subsection{Formal multiple zeta values}

The stuffle regularization map $\operatorname{reg}_*: \h^1_*\simeq\h^0_*[z_1]\to\h^0_*$ is the identity on $\h^0$ and sends $z_1$ to zero. The algebra of \emph{formal multiple zeta values} is
\begin{align}\label{eq:formal-mzv}
 \Zf=\h^0_*\big/\mathrm{EDS}_*,
\end{align}
where $\mathrm{EDS}_*$ is the ideal generated by
\[
 \operatorname{reg}_*(w*v-w\shuffle v),
 \qquad w\in\h^0,\quad v\in\h^1.
\]
For an admissible index we denote the class of $z_{k_1}\cdots z_{k_r}$ by $\zf{k_1,\ldots,k_r}$. There is a natural surjection $\Zf\to\mathcal Z$, and the extended double shuffle conjecture states that it is an isomorphism.

The two polynomial regularizations are available inside the formal algebra itself. Namely, there are unique homomorphisms
\[
 Z^{\mathrm f,*}:\h^1_*\longrightarrow\Zf[T],
 \qquad
 Z^{\mathrm f,\shuffle}:\h^1_\shuffle\longrightarrow\Zf[T]
\]
which extend the quotient map on $\h^0$ and send $z_1$ to~$T$. Put
\begin{align}\label{eq:formal-A}
 A^{\mathrm f}(u)=\exp\left(\sum_{n\geq2}\frac{(-1)^n}{n}\zf{n}u^n\right),
\end{align}
and define the $\Zf$-linear map $\rho^{\mathrm f}:\Zf[T]\to\Zf[T]$ coefficientwise by
\begin{align}\label{eq:formal-rho}
 \rho^{\mathrm f}(e^{Tu})=A^{\mathrm f}(u)e^{Tu}.
\end{align}

\begin{prop}\label{prop:formal-regularization}
One has
\begin{align}\label{eq:formal-reg-theorem}
 Z^{\mathrm f,\shuffle}=\rho^{\mathrm f}\circ Z^{\mathrm f,*}.
\end{align}
Moreover, for every $m\geq1$,
\begin{align}\label{eq:formal-euler}
 \zf{2m}=-\frac{B_{2m}}{2(2m)!}\bigl(-24\zf{2}\bigr)^m.
\end{align}
In particular, $\zf{2m}\in\zf{2}\Zf$.
\end{prop}

\begin{proof}
The proof of the regularization theorem in~\cite[Theorem~2]{IKZ} is an identity in the two polynomial algebras $\h^0_*[z_1]$ and $\h^0_\shuffle[z_1]$. More precisely, after taking the coefficient of $u^n$ in \eqref{eq:formal-rho}, the difference
\[
 Z^{\mathrm f,\shuffle}(w)-\rho^{\mathrm f}
 \bigl(Z^{\mathrm f,*}(w)\bigr)
\]
is the image of a finite $\QQ$-linear combination of the regularized elements
$\operatorname{reg}_*(a*b-a\shuffle b)$ which generate $\mathrm{EDS}_*$. Therefore it is zero in the quotient \eqref{eq:formal-mzv}. This proves \eqref{eq:formal-reg-theorem} without using a numerical realization. Formula \eqref{eq:formal-euler} is likewise a consequence of the extended double shuffle relations, and the formal calculation is given in~\cite[Section~2.2]{BIM}. The last assertion is immediate.
\end{proof}

\section{Formal finite and formal symmetric multiple zeta values}\label{sec:formal-finite}

We now introduce the universal algebra which satisfies the two families of relations common to finite and symmetric multiple zeta values.
This algebra is already considered by Kaneko and Zagier~\cite{KZ},
whose conjecture states that its defining relations give all relations
among finite multiple zeta values. It is studied in detail in the
master's thesis~\cite{Risan}, and a closely related formulation
appears in~\cite[Section~7]{Anzawa}.

\begin{dfn}\label{def:formal-finite}
Let $\mathfrak L\subset\h^1_*$ be the ideal generated by the elements
\begin{align}\label{eq:formal-linear-shuffle}
 w\shuffle v-(-1)^{\wt(w)}R(w)v
\end{align}
for words $w,v\in\h^1$.
The algebra of \emph{formal finite multiple zeta values} is
\[
 \ZAf=\h^1_*/\mathfrak L.
\]
We write $\zaf{k_1,\ldots,k_r}$ for the class of $z_{k_1}\cdots z_{k_r}$ and set $\zaf{\varnothing}=1$.
\end{dfn}

The algebra $\ZAf$ is graded by weight, and we denote its weight-$k$ part by $\ZAfw{k}$. By construction, the stuffle product and the linear shuffle relation hold in $\ZAf$. Specializing \eqref{eq:formal-linear-shuffle} at the empty word gives the reversal relation
\begin{align}\label{eq:formal-reversal}
 \zaf{k_r,\ldots,k_1}=(-1)^{k_1+\cdots+k_r}\zaf{k_1,\ldots,k_r}.
\end{align}

\begin{prop}\label{prop:realizations}
The classical finite and symmetric multiple zeta values define surjective
algebra homomorphisms
\begin{align*}
 \pi_{\mathcal A}:\ZAf&\longrightarrow\ZA,
 &\pi_{\mathcal S}:\ZAf&\longrightarrow
      \mathcal Z/\zeta(2)\mathcal Z,\\
 \zaf{\boldsymbol{k}}&\longmapsto\za{\boldsymbol{k}},
 &\zaf{\boldsymbol{k}}&\longmapsto\zs{\boldsymbol{k}}.
\end{align*}
\end{prop}

\begin{proof}
The stuffle product and \eqref{eq:formal-linear-shuffle} hold for both realizations. The first map is surjective by the definition of $\ZA$. Surjectivity of the second is Yasuda's theorem~\cite{Yasuda}.
\end{proof}

Consequently, every identity in formal finite multiple zeta values proved
below also holds for classical finite
and classical symmetric multiple zeta values.

We will use the following low-depth structure repeatedly in the parity argument.

\begin{prop}\label{prop:low-depth}
For positive integers $a,b$ one has
\begin{enumerate}[(i)]
 \item $\zaf{a}=0$,
 \item if $a+b$ is even, then $\zaf{a,b}=0$,
 \item on setting
 \begin{align}\label{eq:true-single}
  \zzf{1}=0,\qquad \zzf{k}=\frac1k\zaf{k-1,1}\qquad(k\geq2),
 \end{align}
 one has
 \begin{align}\label{eq:formal-double-evaluation}
  \zaf{a,b}=(-1)^a\binom{a+b}{a}\zzf{a+b}.
 \end{align}
\end{enumerate}
\end{prop}

\begin{proof}
We begin with depth one and argue by induction on the weight. The case $k=1$ follows from reversal. For $k=2$, the linear shuffle and stuffle relations for $z_1$ and $z_1$ give, respectively,
\[
 3\zaf{1,1}=0,
 \qquad
 0=\zaf{1}^2=2\zaf{1,1}+\zaf{2},
\]
and hence $\zaf{2}=0$. Let now $k\geq3$. The linear shuffle relation with $w=z_1$ and $v=z_{k-1}$ gives
\begin{align}\label{eq:depth-one-shuffle}
 0=2\zaf{k-1,1}+\sum_{i=2}^{k-2}\zaf{k-i,i}+2\zaf{1,k-1}.
\end{align}
By the induction hypothesis and stuffle,
\[
 \zaf{k-i,i}+\zaf{i,k-i}=-\zaf{k}.
\]
Pairing the terms in \eqref{eq:depth-one-shuffle}, including the middle term separately when $k$ is even, shows that its right-hand side is $-(k+1)\zaf{k}/2$. Thus $\zaf{k}=0$.

We will also use the formal sum formula
\begin{align}\label{eq:formal-sum}
 \sum_{\substack{k_1+\cdots+k_d=k\\k_i\geq1}}
 \zaf{k_1,\ldots,k_d}=0.
\end{align}
Indeed, Hoffman's symmetric-sum identity writes the left-hand side as a polynomial, for the stuffle product, in depth-one values, and these have just been shown to vanish~\cite{Hoffman}. Under the finite realization, \eqref{eq:formal-sum} recovers a special case of the sum formula of Saito--Wakabayashi~\cite{SW}.

If $a+b$ is even, depth-one vanishing, stuffle and reversal give
\[
 0=\zaf{a,b}+\zaf{b,a}=2\zaf{a,b},
\]
which proves part~(ii).

For part~(iii), it is enough to establish the recurrence
\begin{align}\label{eq:double-recurrence}
 (a+1)\zaf{a+1,b}+(b+1)\zaf{a,b+1}=0.
\end{align}
It is immediate from part~(ii) when $a+b+1$ is even. Suppose that this weight is odd. The linear shuffle relation for $z_1$ and $z_az_b$ is
\begin{align}\label{eq:double-linear-shuffle}
0={}&\zaf{1,a,b}
+\sum_{\substack{u+v=a+1\\u,v\geq1}}\zaf{u,v,b}
+\sum_{\substack{u+v=b+1\\u,v\geq1}}\zaf{a,u,v}
+\zaf{a,b,1}.
\end{align}
In addition, expand the following three zero products:
\[
 \zaf{1}\zaf{a,b},\qquad
 \sum_{\substack{u+v=a+1\\u,v\geq1}}\zaf{u}\zaf{v,b},\qquad
 \sum_{\substack{u+v=b+1\\u,v\geq1}}\zaf{u}\zaf{a,v}.
\]
Their stuffle expansions sum to
\begin{align*}
 0={}&\zaf{1,a,b}+\zaf{a,1,b}+\zaf{a,b,1}
 +\zaf{a+1,b}+\zaf{a,b+1}\\
 &+\sum_{u+v=a+1}\bigl(\zaf{u,v,b}+\zaf{v,u,b}+\zaf{v,b,u}\bigr)
 +a\,\zaf{a+1,b}+\sum_{u+v=a+1}\zaf{v,u+b}\\
 &+\sum_{u+v=b+1}\bigl(\zaf{u,a,v}+\zaf{a,u,v}+\zaf{a,v,u}\bigr)
 +\sum_{u+v=b+1}\zaf{u+a,v}+b\,\zaf{a,b+1},
\end{align*}
with all summation variables positive. Subtract twice
\eqref{eq:double-linear-shuffle}. In the difference, the sums
$\sum\zaf{u,v,b}$ and $\sum\zaf{a,u,v}$ cancel after relabeling, the
depth-two contractions combine into
$(a+1)\zaf{a+1,b}+(b+1)\zaf{a,b+1}$ together with the complete sum
$\sum_{u+v=a+b+1}\zaf{u,v}$, which vanishes by \eqref{eq:formal-sum},
and the leftover sums $\sum\zaf{v,b,u}$ and $\sum\zaf{u,a,v}$ cancel in
pairs by reversal in odd weight. This leaves
\begin{align}\label{eq:recurrence-remainder}
 0={}&(a+1)\zaf{a+1,b}+(b+1)\zaf{a,b+1}\\
 &+\zaf{a,1,b}-\zaf{1,a,b}-\zaf{a,b,1}.\notag
\end{align}
The last line vanishes. To see this, set
\[
 T=\zaf{b,a,1}+\zaf{1,b,a}+\zaf{a,1,b}.
\]
Expanding stuffle products with one depth-one factor, we get
\[
 0=2T+\bigl(\zaf{1,a,b}+\zaf{a,b,1}+\zaf{b,1,a}\bigr)
   =\bigl(2-(-1)^{a+b}\bigr)T,
\]
where the last equality is reversal. Hence $T=0$, and in odd total weight this is exactly the last line of \eqref{eq:recurrence-remainder}. This proves \eqref{eq:double-recurrence}. Iterating the recurrence from $\zaf{k-1,1}=k\zzf{k}$ gives \eqref{eq:formal-double-evaluation}.
\end{proof}

The element $\zzf{k}$ can be seen as the formal finite replacement for a single zeta value. Under the finite realization it maps to the Bernoulli class $(B_{p-k}/k)_p$, and it vanishes in even weight.

\subsection{Formal symmetric multiple zeta values}

Let $\epsilon_{\wt}(w)=(-1)^{\wt(w)}w$. If $f,g:\h^1\to B$ are linear maps to a commutative algebra, their deconcatenation convolution is
\[
 (f\star g)(w)=\sum_{uv=w}f(u)g(v).
\]
For $\bullet\in\{*,\shuffle\}$ define
\begin{align}\label{eq:formal-symmetric-map}
 Z_{\mathcal S}^{\mathrm f,\bullet}
 =\bigl(Z^{\mathrm f,\bullet}\circ R\circ\epsilon_{\wt}\bigr)\star Z^{\mathrm f,\bullet}.
\end{align}
Explicitly,
\begin{align}\label{eq:formal-symmetric-explicit}
 Z_{\mathcal S}^{\mathrm f,\bullet}(z_{k_1}\cdots z_{k_r})
 =\sum_{j=0}^r(-1)^{k_1+\cdots+k_j}
 \zeta^{\mathrm f,\bullet,T}(k_j,\ldots,k_1)
 \zeta^{\mathrm f,\bullet,T}(k_{j+1},\ldots,k_r).
\end{align}

We prove the properties of \eqref{eq:formal-symmetric-map} inside the formal algebra. We begin with the product which does not require passing modulo $\zf{2}$.

\begin{lem}\label{lem:formal-symmetric-stuffle}
The map $Z_{\mathcal S}^{\mathrm f,*}:\h^1_*\to\Zf[T]$ is an algebra homomorphism.
\end{lem}

\begin{proof}
The maps $\epsilon_{\wt}$ and $R$ are algebra automorphisms of the stuffle algebra: the first assertion follows from additivity of the weight, and the second follows directly by reversing the recursion \eqref{eq:stuffle}. Hence
$Z^{\mathrm f,*}\circ R\circ\epsilon_{\wt}$ and $Z^{\mathrm f,*}$ are stuffle homomorphisms. The deconcatenation coproduct is compatible with the stuffle product,
\[
 \Delta(w*v)=\Delta(w)*\Delta(v),
\]
where the product on $\h^1\otimes\h^1$ is componentwise. Consequently, the convolution of two stuffle homomorphisms is again a stuffle homomorphism. Applying this observation to \eqref{eq:formal-symmetric-map} proves the lemma.
\end{proof}

For an admissible index $\boldsymbol{k}$, with the empty index allowed, define the leading-one series
\begin{align}\label{eq:leading-one-series}
 G_T^\bullet(\boldsymbol{k};X)
 =\sum_{j\geq0}Z^{\mathrm f,\bullet}(z_1^jz_{\boldsymbol{k}})X^j
 \in\Zf[T][[X]].
\end{align}

\begin{lem}\label{lem:leading-one}
For $\bullet\in\{*,\shuffle\}$ one has
\begin{align}
 G_T^\bullet(\boldsymbol{k};X)
 &=e^{XT}G_0^\bullet(\boldsymbol{k};X),\label{eq:leading-one-T}\\
 G_0^\shuffle(\boldsymbol{k};X)
 &=A^{\mathrm f}(X)G_0^*(\boldsymbol{k};X).\label{eq:leading-one-comparison}
\end{align}
\end{lem}

\begin{proof}
The polynomial decomposition $\h^1_\bullet=\h^0_\bullet[z_1]$ gives the translation identity
\[
 G_{T+U}^\bullet(\boldsymbol{k};X)
 =e^{UX}G_T^\bullet(\boldsymbol{k};X).
\]
Indeed, after comparing coefficients of $U^j$, this is precisely the regularization identity obtained by extracting $j$ leading copies of $z_1$. Equivalently, it follows by induction from the defining shuffle or stuffle recursion. Taking $T=0$ proves \eqref{eq:leading-one-T}. On the other hand, \cref{prop:formal-regularization} and \eqref{eq:leading-one-T} give
\[
 G_T^\shuffle(\boldsymbol{k};X)
 =\rho^{\mathrm f}\bigl(e^{XT}G_0^*(\boldsymbol{k};X)\bigr)
 =A^{\mathrm f}(X)e^{XT}G_0^*(\boldsymbol{k};X).
\]
Setting $T=0$ proves \eqref{eq:leading-one-comparison}.
\end{proof}

\begin{prop}\label{prop:formal-symmetric-regularization}
For every $w\in\h^1$, both $Z_{\mathcal S}^{\mathrm f,*}(w)$ and
$Z_{\mathcal S}^{\mathrm f,\shuffle}(w)$ are independent of $T$ and therefore belong to $\Zf$. Moreover,
\begin{align}\label{eq:formal-reg-congruence}
 Z_{\mathcal S}^{\mathrm f,*}(w)
 \equiv Z_{\mathcal S}^{\mathrm f,\shuffle}(w)
 \pmod{\zf{2}\Zf}.
\end{align}
\end{prop}

\begin{proof}
We group the summands in \eqref{eq:formal-symmetric-explicit} according to the maximal block of entries equal to~$1$ through which the cut passes. Write such a portion of the word as
\[
 z_{r_m}\cdots z_{r_1}z_1^h z_{s_1}\cdots z_{s_n},
\]
where $z_{\boldsymbol r}=z_{r_1}\cdots z_{r_m}$ and
$z_{\boldsymbol s}=z_{s_1}\cdots z_{s_n}$ are admissible, and either of them may be empty. The sum of the contributions from the $h+1$ cuts in this block is the coefficient of $X^h$ in
\begin{align}\label{eq:block-series}
 P^\bullet(\boldsymbol r,\boldsymbol s;X)
 =(-1)^{\wt(\boldsymbol r)}
 G_T^\bullet(\boldsymbol r;-X)G_T^\bullet(\boldsymbol s;X).
\end{align}
By \eqref{eq:leading-one-T}, the two exponential factors in \eqref{eq:block-series} are $e^{-XT}$ and $e^{XT}$, and hence cancel. Thus every group of summands, and therefore the full expression \eqref{eq:formal-symmetric-explicit}, is independent of~$T$.

For the comparison of the two regularizations, \eqref{eq:leading-one-comparison} gives
\begin{align}\label{eq:block-comparison}
 P^\shuffle(\boldsymbol r,\boldsymbol s;X)
 =A^{\mathrm f}(-X)A^{\mathrm f}(X)
 P^*(\boldsymbol r,\boldsymbol s;X).
\end{align}
The odd terms in the logarithm cancel, so that
\begin{align}\label{eq:A-even-factor}
 A^{\mathrm f}(-X)A^{\mathrm f}(X)
 =\exp\left(\sum_{n\geq1}\frac{\zf{2n}}{n}X^{2n}\right)
 \equiv1\pmod{\zf{2}\Zf[[X]]}
\end{align}
by \eqref{eq:formal-euler}. Comparing coefficients of $X^h$ in \eqref{eq:block-comparison} and summing over the blocks proves \eqref{eq:formal-reg-congruence}.
\end{proof}

The remaining relation is combinatorial. We give the operator argument in detail, since this is the point where the linear shuffle relation is transferred to the formal symmetric values.

\begin{lem}\label{lem:cut-shuffle}
Let $B$ be a commutative $\QQ$-algebra and let
$Z:\h^1_\shuffle\to B$ be an algebra homomorphism. Define
\[
 Z_{\mathcal S,Z}=(Z\circ R\circ\epsilon_{\wt})\star Z.
\]
Then, for all words $w,v\in\h^1$,
\begin{align}\label{eq:abstract-linear-shuffle}
 Z_{\mathcal S,Z}(w\shuffle v)
 =(-1)^{\wt(w)}Z_{\mathcal S,Z}(R(w)v).
\end{align}
\end{lem}

\begin{proof}
For a word $z_{k_1}\cdots z_{k_r}$, let
\[
 S_i(z_{k_1}\cdots z_{k_r})
 =(-1)^{k_1+\cdots+k_i}
 (z_{k_i}\cdots z_{k_1})\shuffle
 (z_{k_{i+1}}\cdots z_{k_r}),
 \qquad S=\sum_{i=0}^rS_i.
\]
Since $Z$ is a shuffle homomorphism, the definition of the convolution gives
$Z_{\mathcal S,Z}=Z\circ S$. Under the reversal dictionary
$z_{k_1}\cdots z_{k_r}\mapsto(k_r,\ldots,k_1)$ between our convention and the
increasing-index convention of~\cite{KZ}, and after the relabeling
$i\mapsto r-i$, the operators $S_i$ correspond to the cut operators of Kaneko
and Zagier, and \cite[Proposition~10]{KZ} gives
\[
 S\circ S_i=S\qquad(0\leq i\leq r).
\]
If $d=\dep(w)$, then
$S_d(R(w)v)=(-1)^{\wt(w)}w\shuffle v$. Hence
$S(w\shuffle v)=(-1)^{\wt(w)}S(R(w)v)$, and applying $Z$ proves the claim.
\end{proof}

\begin{prop}\label{prop:formal-symmetric-properties}
The common class $Z_{\mathcal S}^{\mathrm f}(w)\in\Zbarf$ of the two expressions in \cref{prop:formal-symmetric-regularization} satisfies, for all words $w,v\in\h^1$,
\begin{align}\label{eq:formal-symmetric-relations}
 Z_{\mathcal S}^{\mathrm f}(w*v)&=Z_{\mathcal S}^{\mathrm f}(w)Z_{\mathcal S}^{\mathrm f}(v),\\
 Z_{\mathcal S}^{\mathrm f}(w\shuffle v)&=(-1)^{\wt(w)}Z_{\mathcal S}^{\mathrm f}(R(w)v).\notag
\end{align}
\end{prop}

\begin{proof}
The first identity holds already for $Z_{\mathcal S}^{\mathrm f,*}$ by \cref{lem:formal-symmetric-stuffle}. For the second, use the shuffle representative $Z_{\mathcal S}^{\mathrm f,\shuffle}$ and apply \cref{lem:cut-shuffle} to $Z=Z^{\mathrm f,\shuffle}$. The two representatives define the same class by \eqref{eq:formal-reg-congruence}.
\end{proof}

\begin{dfn}\label{def:formal-symmetric}
For an index $\boldsymbol{k}=(k_1,\ldots,k_r)$, its \emph{formal symmetric multiple zeta value} is
\[
 \zsf{k_1,\ldots,k_r}
 =Z_{\mathcal S}^{\mathrm f}(z_{k_1}\cdots z_{k_r})\in\Zbarf.
\]
\end{dfn}

The following depth-two evaluation will be used as a check on the definition and to identify the true single values.

\begin{prop}\label{prop:formal-symmetric-depth-two}
For $a,b\geq1$ and $k=a+b$, one has
\begin{align}\label{eq:formal-symmetric-depth-two}
 \zsf{a,b}=
 \begin{cases}
  0,&k\text{ even},\\
  (-1)^a\binom{k}{a}\zf{k},&k\text{ odd}
 \end{cases}
 \qquad\text{in }\Zbarf.
\end{align}
\end{prop}

\begin{proof}
Set $T=0$, which is allowed by \cref{prop:formal-symmetric-regularization}. From \eqref{eq:formal-symmetric-explicit},
\begin{align}\label{eq:formal-symmetric-depth-two-start}
 \zsf{a,b}
 ={}&\zeta^{\mathrm f,*}(a,b)
 +(-1)^a\zeta^{\mathrm f,*}(a)\zeta^{\mathrm f,*}(b)
 +(-1)^k\zeta^{\mathrm f,*}(b,a)
 \pmod{\zf{2}\Zf}.
\end{align}
If $k$ is even, the stuffle relation turns the right-hand side into
\[
 \bigl(1+(-1)^a\bigr)\zeta^{\mathrm f,*}(a)
 \zeta^{\mathrm f,*}(b)-\zf{k}.
\]
This belongs to $\zf{2}\Zf$: if $a$ is odd the product has coefficient zero, while if $a$ is even both factors have even weight and \eqref{eq:formal-euler} applies.

Suppose that $k$ is odd. The depth-two part of the extended double shuffle relation gives
\begin{align}\label{eq:formal-double-zeta-odd}
 \zeta^{\mathrm f,*}(a,b)\big|_{T=0}
 \equiv\frac12\left((-1)^a\binom{k}{a}-1\right)\zf{k}
 \pmod{\zf{2}\Zf}.
\end{align}
More precisely, for $a\geq2$, the Euler decomposition obtained from the regularized double shuffle relations is
\begin{align*}
 \zf{a,b}={}&(-1)^b
 \sum_{\substack{2\leq j\leq k-1\\j\,\mathrm{even}}}
 \left\{\binom{k-j-1}{a-1}+\binom{k-j-1}{b-1}
       +(-1)^b\delta_{j,a}\right\}\zf{j}\zf{k-j}\\
 &+\frac12\left((-1)^a\binom{k}{a}-1\right)\zf{k}.
\end{align*}
This identity holds in $\Zf$ because it is derived before evaluation from the generators of $\mathrm{EDS}_*$, see also~\cite[Section~2.2]{BIM}. Every term in the sum is divisible by $\zf{2}$ by \eqref{eq:formal-euler}, which proves \eqref{eq:formal-double-zeta-odd} for $a\geq2$. If $a=1$, the stuffle relation at $T=0$,
\[
 0=\zeta^{\mathrm f,*}(1,b)+\zeta^{\mathrm f,*}(b,1)+\zf{k},
\]
reduces the assertion to the already proved case $(b,1)$. Thus \eqref{eq:formal-double-zeta-odd} also covers all regularized boundary cases. Substituting it and its version with $a,b$ interchanged into \eqref{eq:formal-symmetric-depth-two-start} gives
$(-1)^a\binom{k}{a}\zf{k}$, as claimed.
\end{proof}

We can now prove the first main theorem.

\begin{thm}\label{thm:formal-kz}
The map
\begin{align*}
 \varphi:\ZAf&\longrightarrow\Zbarf,\\
 \zaf{\boldsymbol{k}}&\longmapsto\zsf{\boldsymbol{k}}
\end{align*}
is a well-defined surjective homomorphism of graded $\QQ$-algebras.
\end{thm}

\begin{proof}
By \cref{prop:formal-symmetric-properties}, the map
$Z_{\mathcal S}^{\mathrm f}$ is a stuffle homomorphism and annihilates every
generator \eqref{eq:formal-linear-shuffle} of $\mathfrak L$. It therefore
factors through $\ZAf$.

For surjectivity, Yasuda's proof that symmetric multiple zeta values generate
$\mathcal Z/\zeta(2)\mathcal Z$ uses only regularized double shuffle
relations, as he notes in the introduction of~\cite{Yasuda}. The same proof
therefore applies to their universal quotient $\Zf$ and shows that the formal
symmetric values span $\Zbarf$.
\end{proof}

\begin{cor}\label{cor:true-single}
For every $k\geq2$,
\[
 \varphi(\zzf{k})=\zf{k}\pmod{\zf{2}\Zf}.
\]
\end{cor}

\begin{proof}
By \cref{prop:formal-symmetric-depth-two},
$\zsf{k-1,1}=k\zf{k}$ in $\Zbarf$: for odd $k$ this is the second case of \eqref{eq:formal-symmetric-depth-two}, while for even $k$ both sides vanish by \eqref{eq:formal-euler}. Divide by~$k$ and use \eqref{eq:true-single}.
\end{proof}

\begin{conj}\label{conj:formal-kz}
The homomorphism $\varphi$ is an isomorphism.
\end{conj}

This conjecture also appears in~\cite[Conjecture~7.6]{Anzawa}.
To locate the difficulty in this conjecture, note that on the
vector space of indices, the cut-shuffle operator underlying the symmetric
construction has kernel exactly the linear span of the finite linear
shuffle expressions~\cite[Proposition~12]{KZ}. The definition of $\ZAf$,
however, quotients by the \emph{stuffle ideal} generated by this span.
Therefore the unresolved point is not another regularization issue, but the
interaction between multiplication and the cut-shuffle kernel, followed by
the formal multiple-zeta quotient. In particular, a multiplicatively
compatible splitting of the cut-shuffle map would prove
\cref{conj:formal-kz}, whereas the kernel statement alone does
not provide such a splitting.

\section{Parity}\label{sec:parity}

For $d\geq0$, let
\begin{align}\label{eq:depth-filtration}
 \Fil_d\ZAfw{k}
 =\Span_{\QQ}\{\zaf{k_1,\ldots,k_r}:r\leq d,\ k_1+\cdots+k_r=k\},
\end{align}
and let $\Fil_d\ZAw{k}\subset\ZA$ be defined in the same way with
$\zaf{\cdot}$ replaced by $\za{\cdot}$.
The product respects weight and depth filtrations. The antipode gives a first parity reduction which is valid in arbitrary depth.

\begin{thm}\label{thm:weak-parity}
If $k\equiv d\pmod2$, then
\begin{align}\label{eq:weak-parity}
 \Fil_d\ZAfw{k}\subseteq\Fil_{d-1}\ZAfw{k}
 +\sum_{\substack{k_1+k_2=k\\d_1+d_2=d\\d_1,d_2\geq2}}
 \Fil_{d_1}\ZAfw{k_1}\,\Fil_{d_2}\ZAfw{k_2}.
\end{align}
\end{thm}

\begin{proof}
Let $w=z_{a_1}\cdots z_{a_d}$ have weight $k$. Apply the quotient map to \eqref{eq:antipode-identity}. Modulo $\Fil_{d-1}\ZAfw{k}$, contractions in $\Sigma$ disappear, and we obtain
\[
 0\equiv\sum_{i=0}^d(-1)^i
 \zaf{a_i,\ldots,a_1}\zaf{a_{i+1},\ldots,a_d}.
\]
The two end terms are
\[
 \zaf{a_1,\ldots,a_d}+(-1)^d\zaf{a_d,\ldots,a_1}
 =\bigl(1+(-1)^{k+d}\bigr)\zaf{a_1,\ldots,a_d}
\]
by reversal. If $k\equiv d\pmod2$, their sum is twice the original value. The terms with a factor of depth one vanish by \cref{prop:low-depth}, leaving exactly the products in \eqref{eq:weak-parity}.
\end{proof}

In depths three and four the preceding proof gives explicit formulas.

\begin{cor}\label{cor:weak-low-depth}
If $a+b+c$ is odd, then
\begin{align}\label{eq:weak-depth-three}
 \zaf{a,b,c}
 =-\frac12\bigl(\zaf{a+b,c}+\zaf{a,b+c}\bigr).
\end{align}
If $a+b+c+d$ is even, then
\begin{align}
 \zaf{a,b,c,d}
 &=-\frac12\bigl(\zaf{a+b,c,d}+\zaf{a,b+c,d}
 +\zaf{a,b,c+d}+(-1)^{a+b}\zaf{a,b}\zaf{c,d}\bigr).
 \label{eq:weak-depth-four}
\end{align}
\end{cor}

\begin{proof}
For $w=z_az_bz_c$, the terms at the first and second cut in
\eqref{eq:antipode-identity} contain a depth-one factor and vanish.
Expanding the antipode of $w$ therefore gives
\[
 0=\zaf{a,b,c}-\zaf{c,b,a}-\zaf{b+c,a}
   -\zaf{c,a+b}-\zaf{a+b+c}.
\]
The last term vanishes, and reversal in odd total weight gives
\eqref{eq:weak-depth-three}.

For $w=z_az_bz_cz_d$, again the terms at the first and third cuts
vanish. At the middle cut one has
\[
 S_*(z_az_b)*z_cz_d=(z_bz_a+z_{a+b})*z_cz_d,
\]
whose second summand vanishes because $\zaf{a+b}=0$. In the antipode
of the full word, all terms with at least two contractions have depth at
most two and even weight, and hence vanish by \cref{prop:low-depth}.
The remaining identity is
\[\begin{aligned}
0={}&\zaf{a,b,c,d}+\zaf{b,a}\zaf{c,d}+\zaf{d,c,b,a}\\
 &+\zaf{c+d,b,a}+\zaf{d,b+c,a}+\zaf{d,c,a+b}.
\end{aligned}\]
Reversal in even total weight turns the last four terms into the original
depth-four value and the three depth-three terms displayed in
\eqref{eq:weak-depth-four}. Together with
$\zaf{b,a}=(-1)^{a+b}\zaf{a,b}$, this proves the formula.
\end{proof}



The product in \eqref{eq:weak-depth-four} is the only obstruction to strong parity in depth four.

\begin{conj}\label{conj:strong-parity}
If $k\equiv d\pmod2$, then
\[
 \Fil_d\ZAfw{k}\subseteq\Fil_{d-1}\ZAfw{k}.
\]
\end{conj}

For the symmetric realization, the image of this conjecture is known in
every depth by~\cite[Corollary to Theorem~4]{KZ}. Note that the result does not lift
back through the surjection $\pi_{\mathcal S}$, and the corresponding
finite statement in arbitrary depth remains open. A proof in $\ZAf$, on the other hand,
gives both realizations at once.

\begin{thm}
\label{thm:strong-parity-four}
\Cref{conj:strong-parity} holds for $d\leq4$.
\end{thm}

\begin{proof}
For $d=1$, all depth-one values vanish by \cref{prop:low-depth}. For
$d=2$, the relevant weight is even, and every depth-two value vanishes by
the same proposition. For $d=3$, the relevant weight is odd, and
\eqref{eq:weak-depth-three} expresses every depth-three value in depth at
most two. Thus only depth four requires an argument.

Fix an even weight $k$. Write $r=a+b$ and $s=c+d$, so that $r+s=k$.
Formula
\eqref{eq:weak-depth-four} and the depth-two evaluation
\eqref{eq:formal-double-evaluation} give
\begin{align}\label{eq:parity-dictionary}
 \zaf{a,b,c,d}\equiv
 \begin{cases}
  \dfrac12(-1)^{a+c}\binom{r}{a}\binom{s}{c}\zzf r\zzf s,
       &r\text{ odd},\\[3pt]
  0,&r\text{ even},
 \end{cases}
 \pmod{\Fil_3\ZAfw{k}}.
\end{align}
Thus it is enough to prove that every product $\zzf r\zzf s$, where
$r,s\geq3$ are odd and $r+s=k$, belongs to $\Fil_3\ZAfw{k}$.

For odd $r\geq3$, set
\begin{align}\label{eq:odd-f-polynomial}
 f_r(X,Y)=\frac{(Y-X)^r-Y^r+X^r}{XY}
 =\sum_{j=1}^{r-1}(-1)^j\binom rj X^{j-1}Y^{r-1-j}.
\end{align}
Let
\[
 \mathfrak F_k(X,Y,Z,W)=
 \sum_{a+b+c+d=k}\zaf{a,b,c,d}
 X^{a-1}Y^{b-1}Z^{c-1}W^{d-1}.
\]
Then \eqref{eq:parity-dictionary} is equivalent to
\begin{align}\label{eq:F-Q-congruence}
 \mathfrak F_k(X,Y,Z,W)&\equiv\frac12\mathfrak Q_k(X,Y,Z,W)
       \pmod{\Fil_3\ZAfw{k}},\\
 \mathfrak Q_k(X,Y,Z,W)&=
 \sum_{\substack{r+s=k\\r,s\geq3\text{ odd}}}
 \zzf r\zzf s f_r(X,Y)f_s(Z,W),\notag
\end{align}
where the sum is over the ordered pairs $(r,s)$.

We use only the linear shuffle relation with a word of depth one and a
word of depth three. Expanding
\eqref{eq:formal-linear-shuffle} for
$w=z_a$ and $v=z_bz_cz_d$, multiplying by
$(-X)^{a-1}Y^{b-1}Z^{c-1}W^{d-1}$ and summing, gives the exact identity
\begin{align}
0={}&\mathfrak F_k(X,Y,Z,W)+\mathfrak F_k(Y-X,Y,Z,W)\notag\\
 &+\mathfrak F_k(Y-X,Z-X,Z,W)
 +\mathfrak F_k(Y-X,Z-X,W-X,W)\notag\\
 &+\mathfrak F_k(Y-X,Z-X,W-X,-X).
\label{eq:five-term-linear-shuffle}
\end{align}
Note that this gives a direct proof of the generating-series identity,
without using a separate shuffle formula.

Insert \eqref{eq:F-Q-congruence} into
\eqref{eq:five-term-linear-shuffle} and set $X=W=0$. Since
\[
 f_r(0,Y)=-rY^{r-2},\qquad f_r(Y,Y)=0,
 \qquad f_s(Z,0)=sZ^{s-2},
\]
only the first and third substitutions survive. We obtain
\begin{align}
0\equiv{}&\sum_{\substack{r+s=k\\r,s\geq3\text{ odd}}}
 \zzf r\zzf s\,sZ^{s-2}\bigl(f_r(Y,Z)-rY^{r-2}\bigr)\notag\\
={}&\sum_{\substack{r+s=k\\r,s\geq3\text{ odd}}}
 \zzf r\zzf s\,sZ^{s-2}
 \sum_{j=1}^{r-2}(-1)^j\binom rjY^{j-1}Z^{r-1-j}
 \pmod{\Fil_3\ZAfw{k}}.
\label{eq:triangular-generating-series}
\end{align}
The coefficient of $Y^{\alpha-1}Z^{k-\alpha-3}$ yields
\begin{align}\label{eq:triangular-product-system}
 \sum_{\substack{r+s=k,\ r,s\geq3\text{ odd}\\\alpha\leq r-2}}
 s\binom r\alpha\zzf r\zzf s\equiv0
 \pmod{\Fil_3\ZAfw{k}}.
\end{align}
Choose successively $\alpha=r_0-2$ for
$r_0=k-3,k-5,\ldots$ down to the smallest odd $r_0\geq k/2$.
In \eqref{eq:triangular-product-system}, pairs with
$\max(r,s)<r_0$ do not occur, while pairs with
$\max(r,s)>r_0$ have already been treated. The remaining coefficient of
$\zzf{r_0}\zzf{k-r_0}$ is nonzero. Downward induction therefore gives
$\zzf r\zzf s\in\Fil_3\ZAfw{k}$ for all pairs. From
\eqref{eq:parity-dictionary} we get the result.
\end{proof}

\begin{rem}
We note that the original proof of \cref{thm:strong-parity-four} in \cite{Risan} differs from the one presented here. There the theorem was proved by an exhaustive computer-assisted calculation using group actions on generating series. As a result, one also obtains an explicit formula for the depth-four finite multiple zeta values in terms of depth three. We refer to \cite{Risan} for this formula.
\end{rem}

\begin{cor}
\label{cor:explicit-true-single-product}
Let $r,s\geq3$ be odd and put $k=r+s$. Then
\begin{align}\label{eq:four-one-product}
 \zaf{r-1,1,s-1,1}
 =\frac12\zaf{r-1,1}\zaf{s-1,1}
  +\frac12\bigl(\zaf{1,r-1,s}-\zaf{1,s-1,r}\bigr).
\end{align}
Consequently,
\begin{align}\label{eq:four-one-symmetry}
 \zaf{r-1,1,s-1,1}+\zaf{s-1,1,r-1,1}
 &=\zaf{r-1,1}\zaf{s-1,1},\\
 \zaf{r-1,1,s-1,1}-\zaf{s-1,1,r-1,1}
 &=\zaf{1,r-1,s}-\zaf{1,s-1,r}.\notag
\end{align}
\end{cor}

\begin{proof}
Apply \eqref{eq:weak-depth-four} to
$\zaf{r-1,1,s-1,1}$. Since $r$ is odd, it gives
\[
 \zaf{r-1,1,s-1,1}
 =\frac12\zaf{r-1,1}\zaf{s-1,1}
 -\frac12\bigl(
 \zaf{r,s-1,1}+\zaf{r-1,s,1}+\zaf{r-1,1,s}\bigr).
\]
Reversal in the even weight $k$ identifies the first triple in
parentheses with $\zaf{1,s-1,r}$. On the other hand, the stuffle
expansion of $0=\zaf1\zaf{r-1,s}$, whose two depth-two contractions
vanish by \cref{prop:low-depth}, gives
\[
 \zaf{r-1,s,1}+\zaf{r-1,1,s}=-\zaf{1,r-1,s}.
\]
This proves \eqref{eq:four-one-product}. Adding and subtracting it from the
formula with $r$ and $s$ interchanged proves
\eqref{eq:four-one-symmetry}.
\end{proof}

\section{The formal double zeta space}\label{sec:double-zeta}

We recall the formal double zeta space of Gangl, Kaneko and
Zagier~\cite{GKZ}.

\begin{dfn}\label{def:double-zeta}
For $k\geq3$, the \emph{formal double zeta space} $\Dk$ is the $\QQ$-vector space generated by symbols
\[
 Z_{r,s},\quad P_{r,s}\quad(r,s\geq1,\ r+s=k),
 \qquad Z_k,
\]
subject to
\begin{align}\label{eq:formal-double-shuffle}
 P_{r,s}
 &=Z_{r,s}+Z_{s,r}+Z_k,\\
 &=\sum_{a+b=k}\left[\binom{a-1}{r-1}+\binom{a-1}{s-1}\right]Z_{a,b}.\notag
\end{align}
A realization of $\Dk$ in a $\QQ$-vector space $V$ is a linear map $\Dk\to V$, i.e. an assignment satisfying \eqref{eq:formal-double-shuffle}.
\end{dfn}

For even $k$, the quotient relevant to finite multiple zeta values is
\begin{align}\label{eq:reduced-double-zeta}
 \Dtk=\Dk/\mathcal P_k,
 \qquad
 \mathcal P_k=\QQ Z_k+\QQ Z_{1,k-1}
 +\Span_{\QQ}\{P_{r,s}:r,s\text{ even}\}.
\end{align}
We continue to write $Z_{r,s}$ and $P_{r,s}$ for their classes in
$\Dtk$. Since $Z_k=0$ there, the first relation in
\eqref{eq:formal-double-shuffle} becomes
\begin{align}\label{eq:product-in-Dtilde}
 P_{r,s}=Z_{r,s}+Z_{s,r}\qquad(r+s=k).
\end{align}
Here $Z_k$ and the product symbols $P_{r,s}$ with $r,s$ even model terms
divisible by $\zeta(2)$, while $Z_{1,k-1}$ corresponds to the regularized
boundary case.

\begin{prop}\label{prop:dimension-Dtilde}
For every even $k\geq4$,
\begin{align}\label{eq:dimension-Dtilde}
 \dim_{\QQ}\Dtk=\frac{k}{2}-2-\dim_{\QQ}\mathcal S_k.
\end{align}
\end{prop}

\begin{proof}
In even weight the $k/2$ odd--odd symbols form a basis of $\Dk$.
The period-polynomial theorem of~\cite{GKZ} shows that
\[
 \dim_{\QQ}\left(\QQ Z_k+
 \Span_{\QQ}\{P_{r,k-r}:r\text{ even}\}\right)
 =\dim_{\QQ}\mathcal S_k+1.
\]
Here
\[
 Z_k=\frac{2}{k+1}
 \sum_{\substack{2\leq r\leq k-2\\r\ {\rm even}}}P_{r,k-r},
\]
so $Z_k$ already belongs to the span of the even product symbols.
Moreover, $Z_{1,k-1}$ is independent of this subspace: in the
$\kappa$-dependent realization of~\cite[(24)]{GKZ}, its image depends on
$\kappa$, whereas the images of $Z_k$ and the even product symbols do
not. Hence
\[
 \dim_{\QQ}\Dtk=\frac{k}{2}
 -\bigl(\dim_{\QQ}\mathcal S_k+1\bigr)-1,
\]
which is \eqref{eq:dimension-Dtilde}.
\end{proof}

The proof of \cref{prop:dimension-Dtilde} rests on the correspondence,
established in~\cite{GKZ}, between the even period polynomials
\eqref{eq:even-period-polynomials} and relations in $\Dk$. In $\Dtk$
this correspondence takes the following explicit form.

\begin{prop}\label{prop:period-polynomial-relations}
Let $k\geq4$ be even. For $p\in W_k^{\even}$ define $q_{r,s}\in\QQ$ for
$r+s=k$ by
\[
 p(X+Y,Y)=\sum_{r+s=k}\binom{k-2}{r-1}q_{r,s}X^{r-1}Y^{s-1}.
\]
Then
\begin{align}\label{eq:period-polynomial-relation}
 \sum_{\substack{r+s=k\\r,s\geq3\ \odd}}
 \bigl(q_{r,s}-q_{k-1,1}\bigr)Z_{r,s}=0
 \qquad\text{in }\Dtk,
\end{align}
and every linear relation among the classes $Z_{r,s}$ with $r,s\geq3$
odd in $\Dtk$ is of this form. The kernel of the assignment
$p\mapsto\eqref{eq:period-polynomial-relation}$ is the line
$\QQ(X^{k-2}-Y^{k-2})$, so that these relations form a space of
dimension $\dim_{\QQ}\mathcal S_k$.
\end{prop}

\begin{proof}
Write $p(X,Y)=\sum_{r+s=k}\binom{k-2}{r-1}p_{r,s}X^{r-1}Y^{s-1}$. By
\cite[Theorem~3]{GKZ}, one has $q_{r,s}-q_{s,r}=p_{r,s}$ and
\begin{align}\label{eq:GKZ-theorem-three}
 \sum_{\substack{r+s=k\\r,s\ \even}}q_{r,s}Z_{r,s}
 \equiv 3\sum_{\substack{r+s=k\\r,s\ \odd}}q_{r,s}Z_{r,s}
 \pmod{\QQ Z_k}
\end{align}
in $\Dk$. Conversely, an element of $\Dk$, written in the basis
$\{Z_{r,s}:r,s\ \odd\}$, lies in
$\QQ Z_k+\Span_{\QQ}\{P_{r,s}:r,s\ \even\}$ if and only if its
coefficient vector is $(q_{r,s})_{r,s\ \odd}$ for a (unique)
$p\in W_k^{\even}$.

Since $p$ contains only even powers, $p_{r,s}=0$ for even $r$ and $s$,
so that $q_{r,s}=q_{s,r}$ there, and
\[
 2\sum_{\substack{r+s=k\\r,s\ \even}}q_{r,s}Z_{r,s}
 =\sum_{\substack{r+s=k\\r,s\ \even}}q_{r,s}\bigl(P_{r,s}-Z_k\bigr).
\]
In $\Dtk$ the even product symbols and $Z_k$ vanish, and
\eqref{eq:GKZ-theorem-three} becomes
\begin{align}\label{eq:period-open-relation}
 \sum_{\substack{r+s=k\\r,s\ \odd}}q_{r,s}Z_{r,s}=0
 \qquad\text{in }\Dtk.
\end{align}
The summand with $r=1$ vanishes. For the summand with $r=k-1$ we use
the sum formula in $\Dtk$: the relation
\eqref{eq:formal-double-shuffle} for $\{r,s\}=\{1,k-1\}$ reads
$P_{1,k-1}=Z_{k-1,1}+\sum_{a=1}^{k-1}Z_{a,k-a}$, while
$P_{1,k-1}=Z_{1,k-1}+Z_{k-1,1}+Z_k=Z_{k-1,1}$ in $\Dtk$. Hence
$\sum_{a=1}^{k-1}Z_{a,k-a}=0$. Removing the even part, which vanishes
as before, gives
\[
 Z_{k-1,1}=-\sum_{\substack{3\leq a\leq k-3\\a\ \odd}}Z_{a,k-a}.
\]
Substituting this into \eqref{eq:period-open-relation} yields
\eqref{eq:period-polynomial-relation}.

Conversely, suppose that $\sum c_{r,s}Z_{r,s}=0$ in $\Dtk$, the sum
being over odd $r,s\geq3$. By the definition
\eqref{eq:reduced-double-zeta} of $\Dtk$ there is a $\mu\in\QQ$ with
\[
 \sum_{\substack{r+s=k\\r,s\geq3\ \odd}}c_{r,s}Z_{r,s}
 -\mu Z_{1,k-1}
 \in\QQ Z_k+\Span_{\QQ}\{P_{r,s}:r,s\ \even\}
\]
in $\Dk$. By the converse part of \cite[Theorem~3]{GKZ} recalled
above, comparing coefficients in the basis $\{Z_{r,s}:r,s\ \odd\}$
gives a $p\in W_k^{\even}$ with $c_{r,s}=q_{r,s}$ for $r,s\geq3$, with
$q_{k-1,1}=0$ and $q_{1,k-1}=-\mu$. Hence
$c_{r,s}=q_{r,s}-q_{k-1,1}$, as claimed.

Finally, $p=X^{k-2}-Y^{k-2}$ satisfies
$p(X+Y,Y)=(X+Y)^{k-2}-Y^{k-2}$, so that $q_{r,s}=1$ for $r\geq2$ and
$q_{1,k-1}=0$. Thus $p$ lies in the kernel. The classes $Z_{r,s}$ with
$r,s\geq3$ odd span $\Dtk$, since the odd--odd symbols form a basis of
$\Dk$, and $Z_{1,k-1}$ and $Z_{k-1,1}$ have just been expressed in terms
of the remaining ones. The space of relations among these spanning
classes therefore has dimension
$(\frac k2-2)-\dim_{\QQ}\Dtk=\dim_{\QQ}\mathcal S_k
=\dim_{\QQ}W_k^{\even}-1$. By rank--nullity, the kernel is exactly the
line $\QQ(X^{k-2}-Y^{k-2})$.
\end{proof}

By strong parity, $\Fil_4\ZAfw{k}=\Fil_3\ZAfw{k}$ in even weight. We use
depth-four representatives because the map below is stated in terms of
$\zaf{r-1,1,s-1,1}$.

\subsection{A generating family in \texorpdfstring{$\Dtk$}{D-tilde}}

Fix an even integer $k\geq4$ and recall the elements $H_{r,s}$ from
\eqref{eq:H-source-element}. By \eqref{eq:product-in-Dtilde}, they satisfy
\begin{align}\label{eq:H-source-symmetry}
 H_{r,s}+H_{s,r}=rsP_{r,s}.
\end{align}

We first record the polynomial lemma used below. Let $V$ be a
$\QQ$-vector space and let $Z\in V[X,Y]$ be homogeneous of degree
$k-2$. Define
\begin{align}\label{eq:H-source-operator}
 \mathcal L(Z)=(X^2+Y^2)Z+XY(X+Y)(\partial_X+\partial_Y)Z.
\end{align}

\begin{lem}\label{lem:dihedral-odd-coefficients}
Suppose that $k\geq4$ is even and
\begin{align}\label{eq:dihedral-anti-invariance}
 Z(Y,X)=-Z(X,Y),\qquad Z(X,X-Y)=-Z(X,Y).
\end{align}
If
\[
 [X^rY^{k-r}]\mathcal L(Z)=0
 \qquad(3\leq r\leq k-3,\ r\text{ odd}),
\]
then $Z=0$.
\end{lem}

\begin{proof}
The two transformations in \eqref{eq:dihedral-anti-invariance}
generate a dihedral reflection group of order $12$, of type $I_2(6)$,
containing six reflections. Its ring of invariants is
$\QQ[A_2,A_6]$, and the anti-invariant polynomials form the free
module over it generated by the product $J$ of the six linear forms
vanishing on the reflection lines~\cite[Chapter~3]{Humphreys},
where
\[
 J=XY(X-Y)(X+Y)(2X-Y)(X-2Y),\quad
 A_2=X^2-XY+Y^2,\quad A_6=X^2Y^2(X-Y)^2.
\]
A polynomial which is anti-invariant under this group is therefore of
the form
\begin{align}\label{eq:dihedral-factorization}
 Z=J\sum_{\substack{b\geq0\\6b\leq k-8}}
 v_bA_2^{(k-8-6b)/2}A_6^b,
\end{align}
with $v_b\in V$.
For $k=4,6$ the sum is empty, and there is nothing to prove.

Set $m=k-2$, $z(t)=Z(t,1)$ and $w(t)=\mathcal L(Z)(t,1)$. Homogeneity gives
\begin{align}\label{eq:dihedral-one-variable}
 w(t)=\bigl((m+1)t^2+mt+1\bigr)z(t)+(t-t^3)z'(t).
\end{align}
If $z(t)=\sum c_it^i$ and $w(t)=\sum d_jt^j$, then
\begin{align}\label{eq:dihedral-recurrence}
 d_j=(j+1)c_j+mc_{j-1}+(m-j+3)c_{j-2}.
\end{align}
The polynomial $\mathcal L(Z)$ is antisymmetric. Hence the assumption on its odd
coefficients gives
\begin{align}\label{eq:dihedral-boundary}
 w_{\mathrm{odd}}(t)=a(t-t^{m+1}),\qquad a=2z'(0).
\end{align}

Writing \eqref{eq:dihedral-factorization} at $Y=1$ gives
\[
 z(t)=J(t)F\bigl(t^2-t+1,t^2(t-1)^2\bigr),
 \qquad J(t)=J(t,1)=t(t-1)(t+1)(2t-1)(t-2).
\]
Since $J'(0)=J'(1)=-2$, we have $z'(1)=z'(0)$. Equation
\eqref{eq:dihedral-one-variable} now gives
$w'(1)=ma$ and $w'(-1)=0$. Thus
$w_{\mathrm{odd}}'(1)=ma/2$, whereas
\eqref{eq:dihedral-boundary} gives
$w_{\mathrm{odd}}'(1)=-ma$. It follows that $a=0$.

Order the terms in \eqref{eq:dihedral-factorization} by $b$. The
$b$-th term begins with $-2v_bt^{2b+1}$. Once the preceding terms have
vanished, \eqref{eq:dihedral-recurrence} with $j=2b+1$ gives
$-4(b+1)v_b=0$. Induction gives $v_b=0$ for every $b$.
\end{proof}

\begin{prop}\label{prop:H-spans-Dtilde}
For every even $k\geq4$, the elements
\[
 H_{r,s}\qquad(r,s\geq3\text{ odd},\ r+s=k)
\]
span $\Dtk$.
\end{prop}

\begin{proof}
Work in the quotient of $\Dtk$ by the span of the $H_{r,s}$. By
\eqref{eq:H-source-symmetry},
$P_{r,k-r}=0$ for odd $r$ with $3\leq r\leq k-3$.
The same equality for even $r$ follows from the definition of $\Dtk$.
The double shuffle relation for the pair $\{1,k-1\}$, together with
$Z_{1,k-1}=0$, gives
\[
 \sum_{r=1}^{k-1}Z_{r,k-r}=0.
\]
Pairing the terms indexed by $r$ and $k-r$ shows that
$P_{1,k-1}=0$. Thus
\[
 P_{r,k-r}=0\qquad(1\leq r\leq k-1).
\]

Form
\[
 Z(X,Y)=\sum_{r=1}^{k-1}Z_{r,k-r}X^{r-1}Y^{k-r-1}.
\]
These equalities and \eqref{eq:formal-double-shuffle} give
\[
 Z(Y,X)=-Z(X,Y),\qquad Z(X,X-Y)=-Z(X,Y).
\]
Using $P_{a,k-a}=0$, and hence
$Z_{a,k-a}=-Z_{k-a,a}$, in \eqref{eq:H-source-element} we obtain
\[
 2H_{r,s}=-(s+1)Z_{r-1,s+1}-(k-2)Z_{r,s}
 -(r+1)Z_{r+1,s-1}.
\]
The negative of the right-hand side is
$[X^rY^s]\mathcal L(Z)$. All the required coefficients of $\mathcal L(Z)$ therefore vanish,
and \cref{lem:dihedral-odd-coefficients} gives $Z=0$.
\end{proof}

Since there are $k/2-2$ elements $H_{r,s}$,
\cref{prop:dimension-Dtilde,prop:H-spans-Dtilde} also give
\begin{align}\label{eq:H-relation-dimension}
 \dim_{\QQ}\left\{(c_{r,s}):
 \sum_{r+s=k}c_{r,s}H_{r,s}=0\right\}
 =\dim_{\QQ}\mathcal S_k,
\end{align}
where the sum is over odd $r,s\geq3$.

\subsection{The maps \texorpdfstring{$\beta_k$ and $\beta_k^{\mathrm f}$}{beta-k and beta-k-f}}

The elements $H_{r,s}$ span $\Dtk$ by \cref{prop:H-spans-Dtilde}.
Hence, for odd $r,s\geq3$ with $r+s=k$, the rules
\begin{align}\label{eq:direct-beta-rules}
 \beta_k^{\mathrm f}(H_{r,s})
 &=\zaf{r-1,1,s-1,1},\notag\\
 \beta_k(H_{r,s})
 &=\za{r-1,1,s-1,1}
\end{align}
determine at most one linear map in each case, and such a map exists
if and only if the right-hand sides satisfy every linear relation
among the $H_{r,s}$.
If $\mathcal S_k=0$, the elements $H_{r,s}$ form a basis by
\cref{eq:dimension-Dtilde,prop:H-spans-Dtilde}, and both rules define a
map without any further assumption. In the cusp weights this is a
nontrivial condition.

The construction of these maps uses the following two identities.

\begin{conj}\label{conj:beta-identities}
For every even $k\geq4$ and all $a,b,c\geq1$ with $a+b+c=k$, one has
in $\ZAfw{k}$
\begin{align}\label{eq:formal-Hoffman-depth-three}
 \zaf{a,b,c}
={}&(-1)^{b-1}
 \sum_{\substack{n+m=k-1\\n,m\geq1}}
 \binom na\binom mc\,\zaf{n,1,m}\notag\\
&+(-1)^{a+c}\sum_{i=1}^{b-1}
 \binom{a+i}{a}\binom{b+c-i}{c}
 \zzf{a+i}\zzf{b+c-i}.
\end{align}
\end{conj}

The second identity used below, the \emph{harmonic identity}, is a
formal consequence of \cref{conj:beta-identities}.

\begin{prop}\label{prop:harmonic-from-hoffman}
Let $k\geq4$ be even and assume that \cref{conj:beta-identities} holds
in weight $k$. Then, for $2\leq r,s\leq k-2$ with $r+s=k$, one has in
$\ZAfw{k}$
\begin{align}\label{eq:formal-harmonic}
0=
 \sum_{\substack{i+j=r,\ a+b=s\\i,j,a,b\geq1\\i+a\ \mathrm{odd}}}
 \zaf{i,a}\zaf{j,b}
 +2(-1)^r\sum_{\substack{u+v=k-1\\u,v\geq1}}
 \left(\binom vr+\binom vs\right)\zaf{u,1,v}.
\end{align}
The same implication holds for finite multiple zeta values.
\end{prop}

\begin{proof}
All identities are read in weight $k$. Write
\[
 A(X,Y,Z)=\sum_{a+b+c=k}\zaf{a,b,c}X^{a-1}Y^{b-1}Z^{c-1},
 \qquad
 T(U,V)=\sum_{\substack{u+v=k-1\\u,v\geq1}}\zaf{u,1,v}U^{u-1}V^{v-1},
\]
and set $N_n(X,Y)=(X-Y)^n-(-Y)^n$. Multiplying
\eqref{eq:formal-Hoffman-depth-three} by $X^{a-1}Y^{b-1}Z^{c-1}$ and
summing over $a,b,c\geq1$ with $a+b+c=k$ gives
\begin{align}\label{eq:hoffman-summed}
 A(X,Y,Z)={}&\sum_{\substack{n+m=k-1\\n,m\geq1}}\zaf{n,1,m}\,
 \frac{N_n(X,Y)}{X}\,\frac{N_m(Z,Y)}{Z}\notag\\
 &+Y\sum_{\substack{u+v=k\\u,v\ \odd}}
 \zzf u\zzf v\,f_u(X,Y)f_v(Z,Y),
\end{align}
where $f_u$ is the polynomial \eqref{eq:odd-f-polynomial}. Indeed,
the binomial sums over $a$ and $c$ produce the first term by
\[
 \sum_{a=1}^{n}\binom naX^{a-1}(-Y)^{n-a}=\frac{N_n(X,Y)}{X},
\]
and in the second term we used, for odd $u$,
\[
 \sum_{a=1}^{u-1}\binom ua(-X)^aY^{u-a}=XY\,f_u(X,Y),
\]
together with $\zzf u=0$ for even $u$ and $\zzf1=0$.

Expanding the products $0=\zaf a\zaf{b,c}$ by the stuffle product and
using that all depth-two values of even weight vanish gives
$A(X,Y,Z)+A(Y,X,Z)+A(Y,Z,X)=0$, and in particular
\begin{align}\label{eq:cyclic-specialized}
 A(X,Y,X)+2A(Y,X,X)=0.
\end{align}
We substitute \eqref{eq:hoffman-summed} into
\eqref{eq:cyclic-specialized} and multiply by $X^2Y$. By
\cref{lem:restricted-sum} and reversal one has
$\sum_{u+v=k-1}\zaf{u,1,v}=0$, $T(U,V)=T(V,U)$ and
$T(-U,-V)=-T(U,V)$, and moreover $f_v(X,X)=0$ and
$N_m(X,X)=(-1)^{m+1}X^m$. In particular, the product term of
$A(Y,X,X)$ vanishes, so the product terms contribute exactly
$X^2Y^2\sum\zzf u\zzf v f_u(X,Y)f_v(X,Y)$. For the leading-one terms,
expanding $N_nN_m$ and using the three rules gives
\begin{align*}
 Y\sum_{n+m=k-1}\zaf{n,1,m}\,N_n(X,Y)N_m(X,Y)
 &=2Y^2(Y-X)\,T(Y,Y-X),\\
 2\sum_{n+m=k-1}\zaf{n,1,m}\,(-1)^{m+1}N_n(Y,X)X^{m+1}
 &=2X^2(X-Y)\,T(X,X-Y),
\end{align*}
the two left-hand sides being the leading-one parts of
$X^2Y\,A(X,Y,X)$ and of $2X^2Y\,A(Y,X,X)$. Altogether this
yields
\begin{align}\label{eq:harmonic-boxed}
 X^2Y^2\sum_{\substack{u+v=k\\u,v\ \odd}}
 \zzf u\zzf v\,f_u(X,Y)f_v(X,Y)
 &+2Y^2(Y-X)\,T(Y,Y-X)\notag\\
 &+2X^2(X-Y)\,T(X,X-Y)=0.
\end{align}
For $2\leq r\leq k-2$, the coefficient of $X^rY^{k-r}$ in
\eqref{eq:harmonic-boxed} is exactly \eqref{eq:formal-harmonic}, after
evaluating the products $\zaf{i,a}\zaf{j,b}$ by
\eqref{eq:formal-double-evaluation}. The coefficients of
$X^rY^{k-r}$ with $r\in\{0,1,k-1,k\}$ vanish identically by
\cref{lem:restricted-sum} and reversal. Every step used only the
defining relations of $\ZAf$ together with identities proved in
\cref{sec:formal-finite,sec:parity}, so the argument applies verbatim
after the realization $\pi_{\mathcal A}$.
\end{proof}

\begin{prop}\label{prop:finite-hoffman-duality}
\Cref{conj:beta-identities} holds for finite multiple zeta values: for
every even $k\geq4$, the image of
\eqref{eq:formal-Hoffman-depth-three} under
$\pi_{\mathcal A}$ is a valid identity in $\ZA$ for all
$a,b,c\geq1$ with $a+b+c=k$.
\end{prop}

\begin{proof}
Hoffman's duality for finite multiple zeta
values~\cite{HoffmanFinite} states that
$\zeta^\star_{\mathcal A}(\boldsymbol k)
=-\zeta^\star_{\mathcal A}(\boldsymbol k^\vee)$, where
$\zeta^\star_{\mathcal A}$ denotes the star version and
$\boldsymbol k^\vee$ is the Hoffman dual, obtained by writing every
entry as a sum of ones and exchanging all commas and plus signs.
Building on this, Kina proves the two identities
\begin{align*}
 \zeta_{\mathcal F}\bigl((k_1,k_2,k_3)^\vee\bigr)
 &=\zeta^\star_{\mathcal F}(k_1,k_2,k_3)\\
 &\quad-(-1)^{k_1+k_3}\sum_{i=1}^{k_2-1}
 \binom{k_1+i}{k_1}\binom{k_2+k_3-i}{k_3}
 Z_{\mathcal F}(k_1+i)\,Z_{\mathcal F}(k_2+k_3-i),\\
 \zeta_{\mathcal F}\bigl((k_1,k_2,k_3)^\vee\bigr)
 &=(-1)^{k_2-1}
 \sum_{\substack{n_1+1+n_3=k_1+k_2+k_3\\n_1,n_3\geq1}}
 \binom{n_1}{k_1}\binom{n_3}{k_3}\,
 \zeta_{\mathcal F}\bigl((n_1,1,n_3)^\vee\bigr)
\end{align*}
for both the symmetric and the finite multiple zeta values
$\mathcal F\in\{\mathcal S,\mathcal A\}$, together with the
equality of $\zeta_{\mathcal F}\bigl((k_1,1,k_3)^\vee\bigr)$ and
$\zeta^\star_{\mathcal F}(k_1,1,k_3)$, see
\cite[Lemma~4.1, equation~(5.2) and Section~5]{Kina}. Here
$Z_{\mathcal A}(m)=\pi_{\mathcal A}(\zzf m)$ is the Bernoulli class
$(B_{p-m}/m)_p$.
Combine the two identities and expand all star values into finite
multiple zeta values. The contraction terms in these expansions have
depth at most two and even weight $k$, and therefore vanish by
\cref{prop:low-depth}. Hence each star value equals the
corresponding finite multiple zeta value. Translating Kina's index
convention into ours by reversal \eqref{eq:finite-reversal}, with the
substitution $i\mapsto b-i$ in the product sum, then yields exactly
the image of \eqref{eq:formal-Hoffman-depth-three} under
$\pi_{\mathcal A}$.
\end{proof}

\begin{cor}\label{cor:finite-harmonic}
For every even $k\geq4$, the harmonic identity \eqref{eq:formal-harmonic}
holds for finite multiple zeta values.
\end{cor}

\begin{proof}
Combine \cref{prop:finite-hoffman-duality} with
\cref{prop:harmonic-from-hoffman}.
\end{proof}

We do not know whether \cref{conj:beta-identities} follows in
arbitrary weight from the defining relations of $\ZAf$. By exact
rational row reduction in the quotient of the weight-graded piece of
$\h^1_*$ by the stuffle ideal generated by
\eqref{eq:formal-linear-shuffle}, we verified it by computer for every
even weight $k\leq18$, and \eqref{eq:formal-harmonic} then also holds in
these weights by \cref{prop:harmonic-from-hoffman}.

We will use the following restricted sums. They follow directly from
the formal sum formula, reversal, and the product with $\zaf1=0$.

\begin{lem}\label{lem:restricted-sum}
For every even $k\geq4$,
\begin{align}\label{eq:restricted-sum}
 \sum_{a+b=k-1}\zaf{1,a,b}
 =\sum_{a+b=k-1}\zaf{a,1,b}
 =\sum_{a+b=k-1}\zaf{a,b,1}=0,
\end{align}
where the summation variables are positive.
\end{lem}

\begin{proof}
Let $F$ be the sum of all depth-three values of weight $k$, and denote
the three sums in \eqref{eq:restricted-sum} by $L,M,R$. The formal sum
formula gives $F=0$. Summing the depth $1+2$ linear shuffle relation
gives $2F+L+R=0$, while reversal gives $L=R$. Hence $L=R=0$.
Finally, the sum of $0=\zaf1\zaf{a,b}$ over $a+b=k-1$ gives
$L+M+R=0$. The contraction terms have depth two and even weight and
therefore vanish. Thus $M=0$.
\end{proof}

The next lemma gives a well-definedness criterion for
\eqref{eq:direct-beta-rules}.

\begin{lem}\label{lem:direct-beta-certificate}
Assume that \cref{conj:beta-identities} holds in weight $k$. Then the
first rule in \eqref{eq:direct-beta-rules} defines a linear map
\[
 \beta_k^{\mathrm f}:\Dtk\longrightarrow\Fil_4\ZAfw{k}.
\]
Its image is $\Fil_3\ZAfw{k}$.
\end{lem}

\begin{proof}
We verify compatibility by lifting the prescribed values to the standard
source generators. Put
\[
 \mu_k=\frac12
 \sum_{\substack{1\leq n\leq k-1\\n\ \mathrm{odd}}}
 \zzf n\zzf{k-n}.
\]
For $1\leq r\leq k-1$, define
\begin{align}\label{eq:direct-beta-certificate}
 y_r=(-1)^r\sum_{n=1}^{k-2}
 \left(\binom nr-\frac12\binom{k-1}r\right)
 \left(\zaf{n,1,k-n-1}-\zzf n\zzf{k-n}\right)
 -\delta_{r,1}\mu_k.
\end{align}

Assign $Z_k\mapsto0$ and $Z_{r,k-r}\mapsto y_r$ for
$1\leq r\leq k-1$, and abbreviate
$D_n=\zaf{n,1,k-n-1}-\zzf n\zzf{k-n}$.
We first evaluate the shuffle side
$\sum_{a=1}^{k-1}\bigl[\binom{a-1}{r-1}+\binom{a-1}{s-1}\bigr]y_a$ of
\eqref{eq:formal-double-shuffle}, where $s=k-r$. The inversion
identity
\[
 \sum_{a\geq1}(-1)^a\binom{a-1}{r-1}\binom na=
 \begin{cases}
  (-1)^r,&n\geq r,\\
  0,&n<r,
 \end{cases}
\]
applied in \eqref{eq:direct-beta-certificate} gives
\[
 \sum_{a=1}^{k-1}\binom{a-1}{r-1}y_a
 =(-1)^r\sum_{n=r}^{k-2}D_n
 -\frac{(-1)^r}2\sum_{n=1}^{k-2}D_n
 -\delta_{r,1}\mu_k.
\]
By \cref{lem:restricted-sum} and reversal,
\begin{align}\label{eq:boundary-sums}
 \sum_{n=1}^{k-2}D_n=-2\mu_k,
 \qquad
 \sum_{n=r}^{k-2}D_n+\sum_{n=s}^{k-2}D_n=-2\mu_k-\zzf r\zzf s,
\end{align}
so the shuffle side equals
$-(-1)^r\zzf r\zzf s-(\delta_{r,1}+\delta_{s,1})\mu_k$, that is,
$\zzf r\zzf{k-r}$ for $2\leq r\leq k-2$ (both sides vanish for even
$r$) and $-\mu_k$ for $\{r,k-r\}=\{1,k-1\}$. On the other hand, a
Vandermonde convolution
$\sum_i\binom ui\binom v{r-i}=\binom kr$,
followed by \cref{lem:restricted-sum}, rewrites
\eqref{eq:formal-harmonic} as
\begin{align}\label{eq:harmonic-binomial-direct}
 \sum_{n=1}^{k-2}
 \left(\binom nr+\binom ns-\frac12\binom kr\right)D_n
 =-\zzf r\zzf s.
\end{align}
Since
$\binom{k-1}r+\binom{k-1}{s}=\binom kr$, comparing with
\eqref{eq:direct-beta-certificate} gives
$y_r+y_{k-r}=\zzf r\zzf{k-r}$ for
$2\leq r\leq k-2$. At the boundary,
\cref{lem:restricted-sum} and reversal give
\[
 \sum_{n=1}^{k-2}\left(n-\frac{k-1}{2}\right)D_n=-\mu_k,
\]
which together with the first sum in \eqref{eq:boundary-sums},
substituted in \eqref{eq:direct-beta-certificate}, gives
\[
 y_1=0,\qquad y_{k-1}=-\mu_k.
\]
Thus the stuffle and shuffle values agree in every case. Sending each
product symbol to this common value defines a realization of $\Dk$.
By construction $Z_k\mapsto0$, and the boundary calculation gives
$Z_{1,k-1}\mapsto y_1=0$. Each even product symbol $P_{r,k-r}$ maps to
$\zzf r\zzf{k-r}=0$, so the realization descends to $\Dtk$.

Moreover,
\[
 P_{r,k-r}\longmapsto\zzf r\zzf{k-r}
 \qquad(2\leq r\leq k-2).
\]
By
\eqref{eq:four-one-product},
\begin{align}\label{eq:four-one-bridge}
 2\zaf{r-1,1,s-1,1}={}&rs\zzf r\zzf s\notag\\
 &+\zaf{1,r-1,s}-\zaf{1,s-1,r}.
\end{align}
Substituting \eqref{eq:direct-beta-certificate} in
\eqref{eq:H-source-element}, using
$y_r+y_s=\zzf r\zzf s$, and applying
\eqref{eq:formal-Hoffman-depth-three} with
$(a,b,c)=(1,r-1,s)$ and $(1,s-1,r)$ gives
\[
 H_{r,s}\longmapsto\frac12\left(
 rs\zzf r\zzf s+\zaf{1,r-1,s}-\zaf{1,s-1,r}
 \right).
\]
By \eqref{eq:four-one-bridge}, this is
$\zaf{r-1,1,s-1,1}$. Thus the descended map is the map prescribed
in \eqref{eq:direct-beta-rules}.

It remains to determine its image. Add $\mu_k$ back to the first
coordinate in \eqref{eq:direct-beta-certificate}, and denote the
resulting coordinates by $x_r$. Since $x_{k-1}=-\mu_k$, the spans of
the $x_r$ and the $y_r$ agree. For $1\leq r\leq k-2$, put
\[
 U_r=(-1)^rx_r+\binom{k-1}{r}x_{k-1}.
\]
Then
\[
 U_r=\sum_{n=1}^{k-2}\binom nr
 \left(\zaf{n,1,k-n-1}-\zzf n\zzf{k-n}\right),
\]
and binomial inversion gives
\[
 \zaf{n,1,k-n-1}-\zzf n\zzf{k-n}
 =\sum_{r=n}^{k-2}(-1)^{r-n}\binom rnU_r.
\]
Thus every difference on the left belongs to the image. For
$2\leq n\leq k-2$, \eqref{eq:harmonic-binomial-direct} gives
$\zzf n\zzf{k-n}=x_n+x_{k-n}$. Hence
$\zaf{n,1,k-n-1}$ belongs to the image. This is also true for $n=1$,
since $\zzf 1=0$. Finally,
\eqref{eq:formal-Hoffman-depth-three} expresses every depth-three value
in terms of these values and the products $\zzf n\zzf{k-n}$. The image
therefore contains $\Fil_3\ZAfw{k}$. Conversely, the image is spanned
by the values $\zaf{r-1,1,s-1,1}$, which belong to $\Fil_3\ZAfw{k}$ by
\cref{cor:explicit-true-single-product}. The image is therefore
$\Fil_3\ZAfw{k}$.
\end{proof}

\begin{proof}[Proof of \cref{main:double-zeta-surjection}]
Part~(i) is \cref{prop:H-spans-Dtilde}. For part~(ii),
\cref{lem:direct-beta-certificate} shows that the rule in the statement is
well defined and its image is $\Fil_3\ZAfw{k}$. Strong parity gives
$\Fil_3\ZAfw{k}=\Fil_4\ZAfw{k}$. Since the $H_{r,s}$ span the source, the
image is also the span of the $\zaf{r-1,1,s-1,1}$.
\end{proof}

\begin{proof}[Proof of \cref{main:classical-beta}]
By \cref{prop:finite-hoffman-duality}, the identity
\eqref{eq:formal-Hoffman-depth-three} holds for finite multiple zeta
values, and by \cref{cor:finite-harmonic} so does
\eqref{eq:formal-harmonic}. These two identities are the only
consequences of \cref{conj:beta-identities} used in
\cref{lem:direct-beta-certificate}. Hence the construction there and
the proof of \cref{main:double-zeta-surjection} apply verbatim after
the realization $\pi_{\mathcal A}$, with no assumption on $k$. Finally, the $\frac k2-2$ values
$\za{r-1,1,s-1,1}$ span $\Fil_4\ZAw{k}$, and
$\dim_{\QQ}\Fil_4\ZAw{k}\leq\dim_{\QQ}\Dtk
=\frac k2-2-\dim_{\QQ}\mathcal S_k$ by
\cref{prop:dimension-Dtilde}, so these values satisfy at least
$\dim_{\QQ}\mathcal S_k$ linearly independent relations.
\end{proof}

\begin{rem}\label{rem:kina-comparison}
For $2\leq r\leq k-2$, the finite realization of $y_r$ agrees, after
reversal, with the element $Z_{\mathcal A}(k-r,r)$ of Kina, and the
evaluation of the shuffle side above corresponds to his shuffle
relation~\cite{Kina}. The corresponding stuffle relation is left open
in~\cite{Kina}. Here the harmonic identity of
\cref{cor:finite-harmonic} yields the stuffle evaluation
$y_r+y_{k-r}=\zzf r\zzf{k-r}$, which makes the assignment a
realization of $\Dk$.
\end{rem}

Assuming that \cref{conj:beta-identities} holds,
\cref{prop:dimension-Dtilde} shows that the formal part of
\cref{main:double-zeta}, that $\beta_k^{\mathrm f}$ is an isomorphism,
holds if and only if
\begin{align}\label{eq:sharp-depth-four-bound}
 \dim_{\QQ}\Fil_4\ZAfw{k}\geq
 \frac{k}{2}-2-\dim_{\QQ}\mathcal S_k.
\end{align}
The period-polynomial theorem determines the dimension of the source,
but gives no lower bound for the universal formal finite quotient.
Since $\pi_{\mathcal A}$ maps $\Fil_4\ZAfw{k}$ onto $\Fil_4\ZAw{k}$,
the corresponding statement for $\beta_k$ is the same lower bound for
$\dim_{\QQ}\Fil_4\ZAw{k}$, which is stronger.

\subsection{Explicit formulas and period polynomials}
\label{subsec:explicit-period}

The elements $H_{r,s}$ can be written explicitly in terms of the
odd--odd classes. For odd $m$ with $1\leq m\leq k-3$ and $r+s=k$, set,
following~\cite{GKZ},
\begin{align}\label{eq:lambda-prime}
 \lambda'_{m,k-m}(r,s)
 =\sum_{\ell=0}^{r-2}
 \binom{k-2-\ell}{m-1}\binom{r-1}{\ell}B_{k-m-\ell-1},
\end{align}
where $B_\nu$ denotes the $\nu$th Bernoulli number with
$B_1=-\frac12$, and $B_\nu=0$ for $\nu<0$. Since $Z_k=0$ in $\Dtk$,
the identity \cite[(7)]{GKZ} becomes
\begin{align}\label{eq:GKZ-seven-reduced}
 Z_{m+1,k-m-1}
 =-\frac{2}{m}
 \sum_{\substack{r+s=k\\r,s\ \odd}}
 \lambda'_{m,k-m}(r,s)\,Z_{r,s}
 \qquad\text{in }\Dtk,
\end{align}
where $Z_{1,k-1}=0$ and, by the sum formula from the proof of
\cref{prop:period-polynomial-relations},
$Z_{k-1,1}=-\sum_{3\leq a\leq k-3,\ a\ \odd}Z_{a,k-a}$.

\begin{prop}\label{prop:H-in-Z}
Let $k\geq4$ be even and let $r,s\geq3$ be odd with $r+s=k$. Then, in
$\Dtk$,
\begin{align}\label{eq:H-in-Z}
 2H_{r,s}
 ={}&(rs-r+1)Z_{r,s}+(rs+s-1)Z_{s,r}
 +\sum_{\substack{3\leq\ell\leq k-3\\\ell\ \odd}}
 \left(\binom{k-\ell}{s}-\binom{k-\ell}{r}\right)P_{\ell,k-\ell}
 \notag\\
 &+2\sum_{\substack{r'+s'=k\\r',s'\ \odd}}
 \left(\frac{r+1}{r}\,\lambda'_{r,s}(r',s')
      -\frac{s+1}{s}\,\lambda'_{s,r}(r',s')\right)Z_{r',s'},
\end{align}
where $P_{a,b}=Z_{a,b}+Z_{b,a}$.
\end{prop}

\begin{proof}
Apply \eqref{eq:GKZ-seven-reduced} with $m=s$ and with $m=r$ to the two
even--even terms $(s+1)Z_{s+1,r-1}$ and $-(r+1)Z_{r+1,s-1}$ in
\eqref{eq:H-source-element}.
\end{proof}

Substituting $Z_{1,k-1}=0$ and the sum formula for $Z_{k-1,1}$ in
\eqref{eq:H-in-Z} expresses each $H_{r,s}$ as an explicit rational
linear combination
\begin{align}\label{eq:H-in-Z-matrix}
 2H_{r,s}=\sum_{\substack{3\leq r'\leq k-3\\r'\ \odd}}
 A_{r,r'}\,Z_{r',k-r'}
\end{align}
of the classes $Z_{r',k-r'}$ with $r'$ odd and $3\leq r'\leq k-3$. We
write $A_k=(A_{r,r'})$ for the
resulting square matrix, whose rows and columns are indexed by the
odd integers $3\leq r,r'\leq k-3$.

Since the $H_{r,s}$ satisfy relations in $\Dtk$ whenever
$\mathcal S_k\neq0$, the expansion \eqref{eq:H-in-Z-matrix} is not
unique in $\Dtk$. It becomes unique after passing to the quotient
\begin{align}\label{eq:Dbar}
 \Dbk=\Dk/\bigl(\QQ Z_k+\QQ Z_{1,k-1}\bigr),
\end{align}
of which $\Dtk$ is the further quotient by the even product symbols.
As before, we read the right-hand side of \eqref{eq:H-source-element}
in $\Dk$, with $P_{a,b}=Z_{a,b}+Z_{b,a}+Z_k$, so that the $H_{r,s}$
make sense in every quotient of $\Dk$.

\begin{lem}\label{lem:Dbar-basis}
Let $k\geq6$ be even.
\begin{enumerate}[(i)]
\item The classes $Z_{r,k-r}$ with $r$ odd and $3\leq r\leq k-3$ form
a basis of $\Dbk$.
\item The identity \eqref{eq:H-in-Z-matrix} holds in $\Dbk$, and
$A_k$ is the unique matrix with this property.
\end{enumerate}
\end{lem}

\begin{proof}
The elements $Z_k$ and $Z_{1,k-1}$ are linearly independent in $\Dk$,
as in the proof of \cref{prop:dimension-Dtilde}, so
$\dim_{\QQ}\Dbk=\frac k2-2$. The odd--odd classes $Z_{a,k-a}$ with
$a$ odd form a basis of $\Dk$~\cite[Theorem~2]{GKZ}. The relation
\eqref{eq:formal-double-shuffle} for $\{r,s\}=\{1,k-1\}$ gives the sum
formula
\[
 \sum_{a=1}^{k-1}Z_{a,k-a}=Z_{1,k-1}+Z_k
 \qquad\text{in }\Dk,
\]
while the even sum formula of~\cite[Theorem~1]{GKZ} gives
$\sum_{a\ \even}Z_{a,k-a}=\frac34Z_k$. Subtracting, the odd part
$\sum_{a\ \odd}Z_{a,k-a}$ lies in $\QQ Z_k+\QQ Z_{1,k-1}$, so that
\begin{align}\label{eq:odd-sum-Dbar}
 Z_{k-1,1}=-\sum_{\substack{3\leq r\leq k-3\\r\ \odd}}Z_{r,k-r}
 \qquad\text{in }\Dbk.
\end{align}
Hence the $\frac k2-2$ classes $Z_{r,k-r}$ with $r$ odd and
$3\leq r\leq k-3$ span $\Dbk$, and they form a basis by dimension.
This proves part~(i),
and with it the uniqueness in part~(ii). For the identity, note that
the identity \cite[(7)]{GKZ} holds in $\Dk$ up to a rational multiple
of $Z_k$, so that \eqref{eq:GKZ-seven-reduced}, and with it
\eqref{eq:H-in-Z}, hold already in $\Dbk$. Substituting
$Z_{1,k-1}=0$ and \eqref{eq:odd-sum-Dbar} gives
\eqref{eq:H-in-Z-matrix} in $\Dbk$.
\end{proof}

\begin{prop}\label{prop:Dk-basis}
Let $k\geq6$ be even. The following are equivalent:
\begin{enumerate}[(i)]
\item the elements
\[
 H_{r,s}\quad(r,s\geq3\ \odd,\ r+s=k),
 \qquad Z_k,\qquad Z_{1,k-1}
\]
form a basis of $\Dk$,
\item $\det A_k\neq0$.
\end{enumerate}
Both statements hold whenever $\mathcal S_k=0$.
\end{prop}

\begin{proof}
There are $\frac k2$ elements listed in (i), and
$\dim_{\QQ}\Dk=\frac k2$ by \cite[Theorem~2]{GKZ}. Since $Z_k$ and
$Z_{1,k-1}$ are linearly independent and span the kernel of the
projection $\Dk\to\Dbk$, the elements in (i) form a basis of $\Dk$ if
and only if the classes of the $H_{r,s}$ form a basis of $\Dbk$. By
\cref{lem:Dbar-basis}, the latter holds if and only if
$\det A_k\neq0$.

Suppose that $\mathcal S_k=0$. Then
$\dim_{\QQ}\Dtk=\frac k2-2=\dim_{\QQ}\Dbk$ by
\cref{prop:dimension-Dtilde}, so the surjection $\Dbk\to\Dtk$ is an
isomorphism. The classes of the $H_{r,s}$ span $\Dtk$ by
\cref{prop:H-spans-Dtilde}, hence they span $\Dbk$, and therefore
$\det A_k\neq0$.
\end{proof}

\begin{conj}\label{conj:detA}
For every even $k\geq6$ one has $\det A_k\neq0$. Equivalently, the
elements $H_{r,s}$ together with $Z_k$ and $Z_{1,k-1}$ form a basis
of $\Dk$.
\end{conj}

By \cref{prop:Dk-basis}, the content of \cref{conj:detA} lies in the
cusp weights, where it is the statement that every relation of
\cref{prop:period-polynomial-relations} is realized coefficientwise
by \eqref{eq:H-in-Z-matrix}. In each weight it is one explicit determinant condition with binomial
and Bernoulli number entries, and we confirmed it by exact computation
for every even $k\leq100$. Whenever it holds, every
class $Z_{r,s}$ has a unique expression in the elements $H_{r,s}$,
$Z_k$ and $Z_{1,k-1}$. In particular
\[
 Z_{r,k-r}=2\sum_{\substack{3\leq r'\leq k-3\\r'\ \odd}}
 (A_k^{-1})_{r,r'}\,H_{r',k-r'}
 \qquad\text{in }\Dbk,
\]
and expressions in $\Dtk$ are obtained by projection. We do not know
a closed formula for $A_k^{-1}$, a difficulty already encountered
in~\cite[Section~3]{KZ} for a related triangular system. For $p\in W_k^{\even}$, write
\[
 p(X+Y,Y)=\sum_{r+s=k}\binom{k-2}{r-1}q_{r,s}X^{r-1}Y^{s-1}.
\]
For odd $r$ with $3\leq r\leq k-3$, let $\widehat c_r(p)$ be the determinant
of the matrix obtained from $A_k$ by replacing the row indexed by $r$ by
\[
 \bigl(2(q_{r',k-r'}-q_{k-1,1})\bigr)_{r'} ,
\]
where $r'$ runs over the odd integers from $3$ to $k-3$.

\begin{thm}\label{thm:cofactor-relations}
Let $k\geq6$ be even and $p\in W_k^{\even}$. Then
\begin{align}\label{eq:cofactor-relation}
 \sum_{\substack{3\leq r\leq k-3\\r\ \mathrm{odd}}}
 \widehat c_r(p)H_{r,k-r}=0
 \qquad\text{in }\Dtk.
\end{align}
Consequently,
\begin{align}\label{eq:cofactor-finite-relation}
 \sum_{\substack{3\leq r\leq k-3\\r\ \mathrm{odd}}}
 \widehat c_r(p)\za{r-1,1,k-r-1,1}=0
\end{align}
holds unconditionally for finite multiple zeta values. If
\cref{conj:beta-identities} holds in weight $k$, the analogous relation with
$\za{\cdot}$ replaced by $\zaf{\cdot}$ holds in $\ZAf$.

If $\det A_k\neq0$, these relations, as $p$ varies, span the space of all
relations among the $H_{r,s}$, and they vanish identically only for
$p\in\QQ(X^{k-2}-Y^{k-2})$.
\end{thm}

\begin{proof}
All indices below are odd and lie between $3$ and $k-3$. Let $C_{r,r'}$ be
the cofactor of $A_k$ at $(r,r')$. Expansion along the replaced row gives
\[
 \widehat c_r(p)=2\sum_{r'}(q_{r',k-r'}-q_{k-1,1})C_{r,r'}.
\]
The adjugate identity therefore yields
\[
 \sum_r\widehat c_r(p)A_{r,r'}
 =2\det(A_k)(q_{r',k-r'}-q_{k-1,1}).
\]
Using \eqref{eq:H-in-Z-matrix}, we obtain in $\Dbk$
\[
 \sum_r\widehat c_r(p)\,2H_{r,k-r}
 =2\det(A_k)\sum_{r'}
 (q_{r',k-r'}-q_{k-1,1})Z_{r',k-r'}.
\]
The right-hand side vanishes after projection to $\Dtk$ by
\cref{prop:period-polynomial-relations}, proving
\eqref{eq:cofactor-relation}. Applying the unconditional map $\beta_k$
gives \eqref{eq:cofactor-finite-relation}, and applying $\beta_k^{\mathrm f}$
gives the formal relation whenever that map is defined.

If $\det A_k\neq0$, writing
$b_{r'}=q_{r',k-r'}-q_{k-1,1}$, the coefficient row is
\[
 (\widehat c_r(p))_r=2\det(A_k)\,(b_{r'})_{r'}A_k^{-1}.
\]
The final assertions now follow from
\cref{prop:period-polynomial-relations} and
\eqref{eq:H-relation-dimension}.
\end{proof}

\begin{ex}\label{ex:weight-twelve-transfer-relation}
Let $k=12$ and $p=X^2Y^2(X^2-Y^2)^3$. One has
\[
 A_{12}=
 \begin{pmatrix}
  -39&-\frac{23}{3}&-\frac{35}{3}&-\frac{289}{9}\\[2pt]
  -90&-\frac{278}{5}&-\frac{310}{7}&-\frac{1594}{15}\\[2pt]
  90&\frac{628}{5}&\frac{800}{7}&\frac{1594}{15}\\[2pt]
  93&\frac{23}{3}&\frac{35}{3}&\frac{775}{9}
 \end{pmatrix},
 \qquad \det A_{12}=48672.
\]
The determinants in \cref{thm:cofactor-relations} are
\[
 (\widehat c_3(p),\widehat c_5(p),\widehat c_7(p),\widehat c_9(p))
 =(128,72,144,-16).
\]
After dividing by $8$, \eqref{eq:cofactor-finite-relation} becomes
\eqref{eq:intro-weight-twelve-relation}. In the formal algebra the same
relation holds because \cref{conj:beta-identities} has been verified in
weight~$12$.
\end{ex}

\begin{thebibliography}{99}

\bibitem[A]{Anzawa}
T.~Anzawa,
\textit{A Lie algebra associated with adjoint multiple zeta values},
preprint, arXiv:2511.03177 (2025).

\bibitem[BS]{BS}
S.~Baumard and L.~Schneps,
\textit{Period polynomial relations between double zeta values},
Ramanujan J. \textbf{32} (2013), 83--100.

\bibitem[BIM]{BIM}
H.~Bachmann and J.-W.~van Ittersum,
\textit{Formal multiple Eisenstein series and their derivations},
with an appendix by N.~Matthes,
Adv. Math. \textbf{487} (2026), Paper No.~110739.

\bibitem[GKZ]{GKZ}
H.~Gangl, M.~Kaneko and D.~Zagier,
\textit{Double zeta values and modular forms},
in ``Automorphic forms and zeta functions,'' World Sci. Publ., Hackensack, NJ (2006), 71--106.

\bibitem[HST]{HST}
T.~Hara, K.~Sakugawa and K.~Tasaka,
\textit{Symmetric multiple Eisenstein series},
preprint, arXiv:2601.13626 (2026).

\bibitem[H]{Hoffman}
M.~E.~Hoffman,
\textit{Quasi-shuffle products},
J. Algebraic Combin. \textbf{11} (2000), 49--68.

\bibitem[H2]{HoffmanFinite}
M.~E.~Hoffman,
\textit{Quasi-symmetric functions and mod $p$ multiple harmonic sums},
Kyushu J. Math. \textbf{69} (2015), 345--366.

\bibitem[Hu]{Humphreys}
J.~E.~Humphreys,
\textit{Reflection groups and Coxeter groups},
Cambridge Stud. Adv. Math. \textbf{29}, Cambridge Univ. Press, Cambridge, 1990.

\bibitem[HI]{HoffmanIhara}
M.~E.~Hoffman and K.~Ihara,
\textit{Quasi-shuffle products revisited},
J. Algebra \textbf{481} (2017), 293--326.

\bibitem[IKZ]{IKZ}
K.~Ihara, M.~Kaneko and D.~Zagier,
\textit{Derivation and double shuffle relations for multiple zeta values},
Compos. Math. \textbf{142} (2006), 307--338.

\bibitem[K]{Kaneko}
M.~Kaneko,
\textit{An introduction to classical and finite multiple zeta values},
Publ. Math. Besan\c{c}on, Alg\`ebre Th\'eorie Nombres (2019), 103--129.

\bibitem[Ki]{Kina}
K.~Kina,
\textit{An explicit expression for MZVs in terms of symmetric MZVs},
preprint, arXiv:2607.05795 (2026).

\bibitem[KZ]{KZ}
M.~Kaneko and D.~Zagier,
\textit{Finite multiple zeta values},
to appear in the Proceedings of the 17th MSJ-SI Conference on Modular Forms
and Multiple Zeta Values.

\bibitem[P]{Panzer}
E.~Panzer,
\textit{The parity theorem for multiple polylogarithms},
J. Number Theory \textbf{172} (2017), 93--113.

\bibitem[R]{Risan}
Risan,
\textit{Formal finite multiple zeta values and modular forms},
master's thesis, Nagoya University (2025).

\bibitem[SW]{SW}
S.~Saito and N.~Wakabayashi,
\textit{Sum formula for finite multiple zeta values},
J. Math. Soc. Japan \textbf{67} (2015), 1069--1076.

\bibitem[Ta]{Tasaka}
K.~Tasaka,
\textit{Catalogue of modular relations for double zeta values},
RIMS K\^oky\^uroku \textbf{2267} (2023), 100--121.

\bibitem[T]{Tsumura}
H.~Tsumura,
\textit{Combinatorial relations for Euler--Zagier sums},
Acta Arith. \textbf{111} (2004), 27--42.

\bibitem[Y]{Yasuda}
S.~Yasuda,
\textit{Finite real multiple zeta values generate the whole space $\mathcal Z$},
Int. J. Number Theory \textbf{12} (2016), 787--812.

\end{thebibliography}

\end{document}
